\documentclass[11pt,a4paper,oneside]{report}
\usepackage{iftex}
\ifPDFTeX
  \usepackage[utf8]{inputenc}
  \usepackage[T1]{fontenc}
  \usepackage{lmodern}
\else
  \usepackage{fontspec}
  \setmainfont{Latin Modern Roman}
  \setsansfont{Latin Modern Sans}
  \setmonofont{Latin Modern Mono}
\fi
\usepackage[brazil]{babel}
\usepackage{microtype}
\usepackage[left=2.6cm,right=2.3cm,top=2.5cm,bottom=2.4cm,headheight=15pt]{geometry}
\usepackage{xcolor,graphicx,booktabs,tabularx,array,longtable,amsmath}
\usepackage{tikz}
\usetikzlibrary{positioning,arrows.meta}
\usepackage{eso-pic}
\usepackage{titlesec,fancyhdr,enumitem,listings,tcolorbox}
\usepackage{caption}
\usepackage{xurl}
\usepackage{hyperref}
\definecolor{nexus}{HTML}{153650}
\definecolor{accent}{HTML}{168B9A}
\definecolor{pale}{HTML}{F0F5F8}
% Intensidade da marca d'agua: 0 desativa; 0.04 a 0.08 costuma ser discreto.
\newcommand{\nexusWatermarkOpacity}{0.055}
\newcommand{\nexusWatermarkWidth}{0.43\paperwidth}
\hypersetup{colorlinks=true,linkcolor=nexus,urlcolor=accent,pdftitle={NEXUS - Manual técnico},pdfauthor={Equipe NEXUS}}
\setlength{\parindent}{0pt}
\setlength{\parskip}{6pt}
\setlength{\emergencystretch}{2em}
\renewcommand{\arraystretch}{1.25}
\setlength{\tabcolsep}{7pt}
\linespread{1.12}
\titleformat{\chapter}[display]{\color{nexus}\bfseries}{\large CAPÍTULO \thechapter}{5pt}{\Huge}
\titlespacing*{\chapter}{0pt}{0pt}{20pt}
\titleformat{\section}{\Large\bfseries\color{nexus}}{\thesection}{0.7em}{}
\titleformat{\subsection}{\large\bfseries\color{nexus}}{\thesubsection}{0.7em}{}
\pagestyle{fancy}\fancyhf{}
\fancyhead[L]{\small\textcolor{nexus}{NEXUS | Manual técnico}}
\fancyhead[R]{\small Revisão 2.1}
\fancyfoot[L]{\small UFSM | Cachoeira do Sul}\fancyfoot[R]{\thepage}
\renewcommand{\headrulewidth}{0.4pt}
\fancypagestyle{plain}{\fancyhf{}\fancyfoot[L]{\small NEXUS | UFSM}\fancyfoot[R]{\thepage}\renewcommand{\headrulewidth}{0pt}}
\captionsetup{font=small,labelfont=bf,labelsep=endash}
\setlist{nosep,leftmargin=*}
\newtcolorbox{nota}[1]{colback=pale,colframe=accent,boxrule=0.5pt,arc=2mm,title=#1,fonttitle=\bfseries}
\lstset{basicstyle=\ttfamily\small,backgroundcolor=\color{pale},frame=single,rulecolor=\color{accent},numbers=left,numberstyle=\tiny\color{gray},breaklines=true,columns=fullflexible,keepspaces=true,showstringspaces=false,captionpos=b,keywordstyle=\color{nexus}\bfseries,commentstyle=\color{gray}}
\setcounter{tocdepth}{1}

\begin{document}
\begin{titlepage}
\centering
{\small UNIVERSIDADE FEDERAL DE SANTA MARIA\\CAMPUS CACHOEIRA DO SUL\\ENGENHARIA ELÉTRICA\par}
\vspace{2cm}
{\color{nexus}\fontsize{48}{52}\selectfont\bfseries NEXUS\par}
\vspace{0.3cm}{\Large Manual técnico do sistema\par}
\vspace{0.4cm}{\large Instrumentação, telemetria e integração\\do veículo elétrico experimental\par}
\vspace{1cm}\textcolor{accent}{\rule{0.75\textwidth}{1.5pt}}\par
\vspace{1.0cm}
\begin{minipage}{0.83\textwidth}
\centering
\textbf{Equipe discente desta etapa}\par\vspace{0.25cm}
\begin{tabular}{c}
Adylson Hahn\\[0.12cm]
Érick Ramon Biscaglia\\[0.12cm]
Gabriel Muttoni\\[0.12cm]
Gustavo Reckziegel Machado\\[0.12cm]
Théos Schultz Tompsen
\end{tabular}\par
\vspace{0.65cm}
\textbf{Orientador}\par\vspace{0.15cm}
Prof. Dr. Paulo César Vargas Luz\par
\end{minipage}
\vfill
{\small Revisão documental 2.2 --- baseada no repositório em 11/11/2026\\Cachoeira do Sul, RS --- 2026\par}
\end{titlepage}
\pagenumbering{roman}
\chapter*{Controle do documento}
\addcontentsline{toc}{chapter}{Controle do documento}
Este manual descreve o NEXUS a partir dos arquivos de firmware, da documentação do projeto e das decisões de escopo comunicadas pela equipe até 11 de novembro de 2026. Algumas rotinas legadas ainda podem existir no código, embora estejam fora do escopo desta revisão. \textbf{Configuração no código não comprova montagem física ou ensaio concluído.} Onde a placa instalada, a calibração ou um resultado experimental não estão identificados, o texto registra essa limitação de forma explícita. Antes de alterar conexões ou operar o protótipo, confronte a revisão física com os esquemas e as definições do firmware.
\begin{table}[htbp]\centering
\caption{Histórico de revisões do manual}\label{tab:revisoes}
\begin{tabularx}{\textwidth}{@{}llX@{}}\toprule
Revisão & Data & Conteúdo\\\midrule
0.1 & 10/11/2026 & Modelo inicial e estrutura editorial do arquivo recebido.\\
1.0 & 10/11/2026 & Redação técnica baseada no repositório: arquitetura, ligações, aquisição, protocolos, operação, atualização e pontos de validação.\\
2.0 & 10/11/2026 & Aprofundamento de eletrônica, Wi-Fi, comandos, memória, Google, OTA, rollback, falhas, OLED e procedimentos de bancada.\\
2.1 & 10/11/2026 & Revisão de arquitetura, sensores, firmware e tabelas; retirada da leitura analógica de tensão e do áudio; diagramas, ligações e marca d'água.\\
2.2 & 11/11/2026 & Acesso HTTP, redes abertas e varredura Wi-Fi, leitura INA228 em Operação e sensores, terminal e logs, gráfico elétrico opcional e ajuste dos diagramas.\\\bottomrule
\end{tabularx}\end{table}
\begin{nota}{Leitura correta dos estados}
\textbf{Implementado} significa presente no código analisado; \textbf{configurado} é um valor definido no firmware; \textbf{a verificar} requer confronto com a montagem; \textbf{não validado} significa que não há registro de ensaio suficiente no material consultado. Telemetria simulada serve para demonstração e nunca deve ser confundida com medição.
\end{nota}
\tableofcontents
\listoffigures\addcontentsline{toc}{chapter}{Lista de figuras}
\listoftables\addcontentsline{toc}{chapter}{Lista de tabelas}
\chapter*{Siglas e convenções}\addcontentsline{toc}{chapter}{Siglas e convenções}
\small\begin{longtable}{@{}p{2.2cm}p{13.4cm}@{}}\toprule
Sigla & Significado e uso no NEXUS\\\midrule
ACK & Confirmação de recebimento de um quadro UART; não representa resultado final de uma operação longa.\\
API HTTP & Conjunto de rotas de rede que o navegador consulta nos satélites; distinto dos quadros UART e das funções internas do firmware.\\
AP & Ponto de acesso Wi-Fi criado por um satélite.\\
CRC & Verificação de redundância cíclica; detecta corrupção de quadros, sem autenticar o emissor.\\
DAC & Conversor digital-analógico; produz uma tensão a partir de um valor digital.\\
CSV & Valores separados por vírgulas; formato do histórico de telemetria no microSD.\\
DHCP & Serviço que atribui endereços IP aos clientes do AP.\\
ESP-IDF & Ambiente de desenvolvimento da Espressif usado pelos três firmwares.\\
GPIO & Pino de entrada ou saída de propósito geral.\\
INA228 & Circuito integrado de monitoramento de corrente, tensão e potência da Texas Instruments.\\
HTTP/HTTPS & Protocolo das páginas e APIs; HTTPS acrescenta proteção TLS.\\
I\textsuperscript{2}C & Barramento serial compartilhado pelo MPU6050 e pela tela OLED.\\
LDO & Regulador linear de baixa queda de tensão; dissipa a diferença de tensão em calor.\\
NVS & Armazenamento persistente de pequenos valores na flash.\\
OTA & Atualização de firmware pela rede, sem cabo após a primeira gravação.\\
PCNT & Contador de pulsos por hardware do ESP32-S3, usado nos encoders.\\
PSRAM & Memória RAM externa volátil do módulo mestre.\\
PWM & Modulação por largura de pulso, usada para gerar tons no buzzer.\\
RTC & Domínio de relógio de tempo real do ESP; sua memória pode sobreviver a certos reinícios, mas não a corte de energia.\\
RSSI & Indicador de potência recebida no enlace Wi-Fi, em dBm.\\
RTT & Tempo de ida e volta entre um envio UART e sua confirmação.\\
SD/SPI & Cartão Secure Digital e barramento serial usado para acessá-lo.\\
SHA-256 & Resumo criptográfico de 256 bits usado para conferir imagens OTA.\\
TLS & Proteção criptográfica da comunicação HTTPS.\\
UART & Interface serial assíncrona, com TX e RX entre mestre e satélites.\\
VBUS & Pino/medição de tensão do barramento do banco de baterias no INA228, conforme o ponto de ligação escolhido.\\
VSHUNT & Tensão diferencial entre os terminais de medição do shunt usada para calcular corrente.\\\bottomrule
\end{longtable}\normalsize
\par\medskip Grandezas são indicadas com unidade: tensão em V, corrente em A, potência em W, velocidade em km/h, aceleração em m/s\textsuperscript{2} e tempo em ms ou s. Frequência de aquisição, publicação, registro e atualização do navegador são processos distintos.
\clearpage\pagenumbering{arabic}
\AddToShipoutPictureBG{%
  \begin{tikzpicture}[remember picture,overlay]
    \fill[white] (current page.south west) rectangle (current page.north east);
    \node[opacity=\nexusWatermarkOpacity,blend mode=difference]
      at (current page.center)
      {\includegraphics[width=\nexusWatermarkWidth]{imagens/logo_carro_eletrico.png}};
  \end{tikzpicture}%
}
\chapter{Introdução}
\section{O projeto e a finalidade do NEXUS}
O veículo elétrico experimental da Universidade Federal de Santa Maria, campus Cachoeira do Sul, reúne energia, mecânica, eletrônica e programação numa plataforma de ensino e pesquisa. O NEXUS organiza a instrumentação e a comunicação necessárias para acompanhar o veículo em ensaios: obtém sinais, transforma-os em dados com unidade e indicação de validade, distribui a telemetria e preserva registros para análise. O nome expressa a ligação entre os subsistemas embarcados e as pessoas que os utilizam.

Instrumentação abrange o sensor, sua alimentação e ligação elétrica, a aquisição, o processamento e a calibração. Telemetria é o transporte dos resultados até as interfaces e os arquivos. O sistema oferece um painel técnico à equipe e um painel público aos visitantes. Há também pedidos de controle em desenvolvimento; na revisão documentada, o firmware apenas registra esses pedidos e \textbf{não aciona fisicamente motor, freio ou controlador}.

\section{Objetivo e público}
Este manual reúne a arquitetura, as ligações previstas no firmware, o tratamento dos dados, os serviços embarcados, a operação e os procedimentos de verificação. Destina-se a quem monta e confere o chicote, desenvolve os programas, acompanha ensaios ou interpreta os registros. O projeto do circuito de tração, os dispositivos físicos de proteção e a montagem mecânica completa precisam de seus próprios desenhos, cálculos e ensaios; esta descrição da telemetria não os substitui.

\section{O que está disponível nesta revisão}
Há três programas independentes para ESP32-S3: Mestre, Equipe e Visitantes. O Mestre integra aquisição de movimento, cálculo, histórico, comunicação com os satélites e serviços de rede usados conforme a necessidade. Equipe e Visitantes recebem telemetria por enlaces UART independentes e oferecem páginas locais por Wi-Fi, com níveis diferentes de acesso. O retrato é publicado aos satélites a cada 250\,ms em operação normal; o histórico seleciona uma amostra por segundo quando o microSD está disponível.

O repositório contém mais de uma revisão de circuito, mas não identifica de modo inequívoco qual placa está instalada no veículo. A geometria das rodas e a relação mecânica dos encoders também exigem medição no conjunto real. O modelo de telemetria reserva campos para energia, pedais, localização e outras grandezas; a presença de um campo no pacote ou no painel \textbf{não comprova uma medição física}. O driver do INA228 já foi integrado ao Mestre para leitura elétrica, com shunt configurado de 1\,m$\Omega$ e faixa prevista de 75\,A. A montagem, o sentido da corrente e a calibração ainda exigem conferência na placa real antes de interpretar os valores como ensaio validado.

\begin{nota}{Dados simulados}
O simulador embarcado ajuda a verificar quadros e páginas sem depender de todos os sensores. A configuração padrão usa a origem real; quando a simulação é ativada durante uma sessão, sua origem é marcada no protocolo e deve ser identificada na interface. A partida seguinte restaura o padrão compilado.
\end{nota}

\section{Como usar o manual}
Comece pelo capítulo~\ref{cap:arquitetura} para localizar os três ESPs e seguir o caminho de uma amostra. Use o capítulo~\ref{cap:hardware} para conferir conectores, GPIOs, níveis lógicos e sensores \emph{antes} de ligar a placa; em seguida, consulte o capítulo~\ref{cap:eletronica} para alimentação, barramentos, INA228 e circuito de habilitação do driver. As tabelas de pinagem descrevem o software e precisam ser comparadas com a revisão física montada.

Para alterar ou investigar o funcionamento, siga o capítulo~\ref{cap:firmware} e os capítulos dedicados a comandos, Wi-Fi, memória, atualização e diagnóstico. Eles diferenciam o que já funciona, o que está apenas previsto e o que ainda requer ensaio. Na operação, use o capítulo~\ref{cap:operacao} como sequência de preparação e registro de resultados. Ao encontrar um valor inesperado, verifique sua origem, unidade, validade e versão do firmware antes de compará-lo com uma referência. O apêndice identifica as fontes e os dados mínimos que devem acompanhar cada ensaio.

\chapter{Arquitetura do sistema}\label{cap:arquitetura}
\section{Divisão em três unidades}
O Mestre coordena os serviços embarcados: recebe sinais, forma o retrato do veículo, distribui telemetria, registra histórico e executa as operações que precisam de rede externa. A unidade Equipe apresenta o estado técnico e encaminha solicitações de manutenção e configuração. A unidade Visitantes oferece uma visualização pública de leitura. Separar essas funções permite que cada interface evolua e atenda seus usuários sem comandar diretamente a cadência de aquisição do Mestre. Cada satélite usa seu próprio enlace UART com ele.

\begin{figure}[htbp]\centering
\begin{tikzpicture}[>={Latex}, every node/.style={draw=nexus,rounded corners,fill=pale,align=center,minimum height=10mm,font=\small}]
\node[text width=3.0cm] (entradas) at (0,2.3) {Sensores e sinais\\do veículo};
\node[text width=3.4cm] (mestre) at (0,0) {ESP32-S3 Mestre\\processamento e serviços};
\node[text width=2.5cm] (sd) at (-4.8,0) {microSD\\histórico e logs};
\node[text width=2.9cm] (equipe) at (-3.3,-2.4) {ESP32-S3 Equipe\\painel técnico};
\node[text width=2.9cm] (visitantes) at (3.3,-2.4) {ESP32-S3 Visitantes\\painel público};
\node[text width=2.9cm] (operador) at (-3.3,-4.25) {Navegador da equipe};
\node[text width=2.9cm] (publico) at (3.3,-4.25) {Navegador visitante};
\draw[->,thick] (entradas.south)--(mestre.north);
\draw[<->,thick] (sd.east)--node[draw=none,fill=white,above,font=\scriptsize,inner sep=2pt]{SPI2}(mestre.west);
\draw[<->,thick] (mestre.south west)--node[draw=none,fill=white,above,sloped,font=\scriptsize,inner sep=2pt]{UART1}(equipe.north east);
\draw[<->,thick] (mestre.south east)--node[draw=none,fill=white,above,sloped,font=\scriptsize,inner sep=2pt]{UART2}(visitantes.north west);
\draw[<->,thick] (equipe.south)--node[draw=none,fill=white,right,font=\scriptsize,inner sep=2pt]{Wi-Fi/HTTP}(operador.north);
\draw[<->,thick] (visitantes.south)--node[draw=none,fill=white,right,font=\scriptsize,inner sep=2pt]{Wi-Fi/HTTP}(publico.north);
\end{tikzpicture}
\caption{Fluxo entre unidades. Os dois enlaces UART e os dois pontos de acesso Wi-Fi são independentes.}\label{fig:arquitetura}
\end{figure}

\begin{table}[htbp]\centering\small
\caption{Responsabilidade principal de cada unidade}\label{tab:unidades}
\begin{tabularx}{\textwidth}{@{}p{2.7cm}X@{}}\toprule
Unidade & Responsabilidade no sistema\\\midrule
Mestre & Coordena a aquisição e o processamento, mantém o estado do sistema, publica telemetria, registra os dados de bordo e conduz operações externas, como atualização e envio de arquivos.\\
Equipe & Mantém a interface técnica local, mostra diagnósticos e faz a ponte entre as solicitações do operador e o Mestre.\\
Visitantes & Mantém a interface pública local para acompanhar a telemetria, sem oferecer comandos operacionais.\\\bottomrule
\end{tabularx}\end{table}

\section{Da aquisição à apresentação}
Cada fonte de dados possui sua própria forma de aquisição. O Mestre combina os resultados disponíveis em um retrato \texttt{DadosTelemetriaVeiculo}, preservando origem e indicadores de validade. O serviço de telemetria converte valores para as escalas do protocolo, atribui sequência e tempo desde a partida e envia cópias às UARTs. Cada satélite verifica formato, versão, tamanho, CRC e plausibilidade antes de disponibilizar a amostra ao seu servidor HTTP. O ACK devolvido ao Mestre permite medir prazo de resposta e perdas do enlace.

Em paralelo, o serviço de histórico seleciona uma amostra por segundo e a entrega a uma tarefa que escreve no microSD. Essa tarefa não dita o ritmo de aquisição: uma escrita demorada é absorvida pela fila até seu limite, e eventuais perdas são contabilizadas. O navegador consulta o satélite a seu próprio ritmo; fechar ou recarregar a página não apaga o CSV de bordo. As páginas e as rotas HTTP são detalhadas no capítulo~\ref{cap:wifi-paineis}.

\section{Cadências dos serviços}
Uma frequência de amostragem, uma frequência de publicação e uma frequência de atualização da página medem atividades diferentes. Não há uma frequência única para todos os sensores: alguns são atendidos por tarefa periódica e os pulsos dos encoders são acumulados pelo periférico PCNT quando ocorrem. A tabela apresenta os ritmos dos \emph{serviços}, sem confundi-los com a velocidade elétrica dos barramentos.

\begin{table}[htbp]\centering\small
\caption{Periodicidade dos serviços na configuração consultada}\label{tab:ritmos}
\begin{tabularx}{\textwidth}{@{}p{3.3cm}p{3.2cm}X@{}}\toprule
Serviço & Periodicidade & O que o intervalo determina\\\midrule
Aquisição & Própria de cada fonte & Leituras periódicas ou eventos são processados para formar o próximo retrato.\\
Telemetria normal & 250\,ms (4\,Hz) & Novo retrato enviado a cada satélite pela sua UART.\\
Histórico no microSD & 1\,s (1\,Hz) & Seleção de uma amostra para o CSV, se o cartão estiver disponível.\\
Consulta do navegador & Equipe: 250--1000\,ms; Visitantes: 500\,ms & Intervalo solicitado pela página; rede e servidor podem ampliar o prazo efetivo.\\
Diagnóstico do Mestre & 1\,s & Atualização do estado de saúde e das causas ativas.\\
Telemetria durante OTA & 2\,s por padrão & Reserva mais tempo da UART para blocos da imagem; ajustável de 500\,ms a 10\,s.\\\bottomrule
\end{tabularx}\end{table}

I\textsuperscript{2}C reúne os periféricos compatíveis em um barramento compartilhado, SPI2 liga o microSD e PCNT conta bordas dos encoders sem exigir leitura contínua de pinos pela CPU. A rede externa do Mestre é usada em sessões de OTA e envio ao Google; o fluxo local de aquisição, UART e painéis não depende desse roteador.

\section{Falhas, degradação e indicação ao operador}
\emph{Operação degradada} significa que o Mestre continua executando as funções essenciais, porém perdeu uma função ou reduziu sua qualidade: por exemplo, sem microSD há telemetria ao vivo, mas não CSV persistente; sem tela ainda há UART e painel; sem roteador externo, OTA e Google ficam indisponíveis, enquanto a comunicação local continua. Já \emph{falha} é o estado de saúde usado para condições críticas, especialmente pouca memória ou margem de pilha insuficiente. A aplicação pode reiniciar se a partida de um serviço essencial não puder ser concluída. Nem toda ausência de módulo opcional tem a mesma gravidade.

O diagnóstico acompanha resultados de sensores, erros I\textsuperscript{2}C e UART, falta de resposta dos satélites, atrasos de tarefas, memória e pilhas, intensidade de sinal da rede externa, persistência e falhas OTA. Ele combina contagem, persistência e janelas temporais para distinguir um erro isolado de um problema recorrente. O estado pode ser \emph{saudável}, \emph{atenção}, \emph{falha} ou \emph{recuperando}; a causa e os limiares de cada caso estão no capítulo~\ref{cap:diagnostico}.

O operador vê o estado resumido nos LEDs e no buzzer, o motivo e o progresso na OLED, e os detalhes nos relatórios seriais e no painel técnico da Equipe. A telemetria informa origem, validade e sequência dos dados que chegam aos satélites; os diagnósticos trazem contadores e causas para investigar perda de função. Verde sinaliza condição saudável, amarelo atenção, vermelho falha crítica e azul marca fases como partida ou atualização, conforme o padrão de sinalização. Avisos sonoros marcam transições relevantes. Assim, a ausência de um registro no SD, por exemplo, pode ser identificada sem atribuir automaticamente defeito aos sensores de movimento.

\chapter{Hardware, sensores e ligações}\label{cap:hardware}
\section{Placas ESP32-S3 e restrições de pinos}
O Mestre usa um ESP32-S3-WROOM-1-N16R8, configurado para 16\,MiB de flash e 8\,MiB de PSRAM Octal. Os satélites usam placas ESP32-S3 SuperMini com flash configurada para 4\,MiB. Os três processadores possuem dois núcleos Xtensa LX7; o chip admite até 240\,MHz, e os três projetos fixam 160\,MHz nos respectivos arquivos \texttt{sdkconfig}. Clock do processador, capacidade da flash e memória externa são propriedades diferentes: aumentar a primeira não cria espaço para programa ou histórico.

\begin{table}[htbp]\centering\small
\caption{Placas e recursos usados pelos três programas}\label{tab:modelos-esp}
\begin{tabularx}{\textwidth}{@{}p{2.5cm}p{4.2cm}X@{}}\toprule
Unidade & Placa/módulo previsto & Recursos e conferência necessária\\\midrule
Mestre & ESP32-S3-WROOM-1-N16R8 & Dois núcleos a 160\,MHz; flash de 16\,MiB e PSRAM Octal de 8\,MiB. Conferir o módulo soldado e a entrada de alimentação da placa.\\
Equipe e Visitantes & ESP32-S3 SuperMini & Dois núcleos a 160\,MHz e flash configurada para 4\,MiB em cada placa. Confirmar fabricante, pinagem e capacidade física da variante adquirida.\\\bottomrule
\end{tabularx}\end{table}

O repositório contém mais de uma revisão de circuito e não identifica de forma inequívoca a PCB montada. A tabela de GPIO abaixo descreve \textbf{a atribuição no firmware}, que deve ser comparada com conectores e trilhas da placa real. Antes de energizar, conferir polaridade e tensão dos trilhos, continuidade do GND, proteções e ausência de curto. Não aplicar 12 ou 24\,V diretamente aos GPIOs: UART, I\textsuperscript{2}C, entradas dos encoders e demais sinais lógicos do ESP trabalham em 3,3\,V.

\begin{table}[htbp]\centering\small
\caption{GPIOs que exigem reserva ou cuidado no ESP32-S3}\label{tab:pinos-especiais}
\begin{tabularx}{\textwidth}{@{}p{2.6cm}X@{}}\toprule
GPIO & Motivo e decisão de projeto\\\midrule
26--32 & Normalmente ligados à flash e à PSRAM pelo módulo; não destiná-los ao chicote.\\
33--37 no N16R8 & Ocupados pela interface da PSRAM Octal do Mestre; não destiná-los a periféricos externos. Em outra variante de placa, consultar seu esquema.\\
0, 3, 45 e 46 & Pinos de configuração de partida (\emph{strapping}); níveis externos indevidos durante reset podem alterar o modo de boot. Usar somente após analisar a placa e os resistores de polarização.\\
19 e 20 & Usados pelo USB nativo/Serial-JTAG; reaproveitá-los como GPIO pode impedir gravação ou diagnóstico por essa porta.\\\bottomrule
\end{tabularx}\end{table}

Essas restrições vêm da \href{https://docs.espressif.com/projects/esp-idf/en/stable/esp32s3/api-reference/peripherals/gpio.html}{documentação de GPIO do ESP32-S3}; pinos de \emph{strapping} e USB não são universalmente proibidos, mas devem preservar a função de partida e manutenção exigida na placa. O clock máximo e os dois núcleos constam da \href{https://www.espressif.com/sites/default/files/documentation/esp32-s3-wroom-1_wroom-1u_datasheet_en.pdf}{folha de dados do módulo}.

\section{Mapa de pinos do mestre}
\begin{table}[htbp]\centering\small
\caption{Interfaces do mestre definidas no firmware}\label{tab:pinos-mestre}
\begin{tabularx}{\textwidth}{@{}p{3.0cm}p{2.4cm}X@{}}\toprule
Interface & GPIO do Mestre & Ligação e função\\\midrule
Encoder esquerdo & 5 (A), 6 (B) & Entradas de pulsos do lado esquerdo; cada canal recebe pull-up externo a 3,3\,V.\\
Encoder direito & 7 (A), 21 (B) & Entradas independentes do lado direito, com a mesma referência GND e pull-ups.\\
I\textsuperscript{2}C & 8 (SDA), 9 (SCL) & Barramento compartilhado por MPU6050, OLED e INA228, configurado a 400\,kHz; confirmar no esquema e na montagem o conector efetivamente usado pelo INA228.\\
LEDs & 10, 11, 12, 13 & Verde, amarelo, vermelho e azul, respectivamente; saídas ativas em alto e resistor em série por LED.\\
Buzzer passivo & 14 & Saída de tons por PWM; dimensionar o estágio de acionamento para o buzzer instalado.\\
UART Visitantes & 15 (TX), 16 (RX) & Canal dedicado ao satélite público; cruzar TX com RX do outro ESP.\\
UART Equipe & 17 (TX), 18 (RX) & Canal dedicado ao satélite técnico; não unir as duas UARTs.\\
microSD SPI2 & 38 (CLK), 39 (MOSI), 40 (MISO), 41 (CS) & Barramento do cartão, configurado a 10\,MHz; conferir alimentação e nível lógico do módulo.\\\bottomrule
\end{tabularx}\end{table}

\section{Encoders e medição do movimento}
O projeto prevê um encoder incremental E38S6G5-600B-G24N em cada lado do veículo. Cada unidade fornece dois sinais de pulsos, A e B, defasados no tempo. O modelo documentado tem 600 pulsos por volta do próprio eixo. A defasagem permitiria deduzir o sentido de rotação com decodificação em quadratura; nesta revisão, o firmware conta as bordas dos dois canais para melhorar a resolução, mas \textbf{não calcula o sentido}. Por isso, recuar a roda também aumenta o percurso acumulado.

As saídas são NPN de \emph{coletor aberto}. O transistor interno do encoder apenas liga a saída ao GND quando conduz; desligado, ele deixa a saída eletricamente aberta. Sem um resistor de \emph{pull-up}, esse estado aberto ficaria indefinido. Cada canal A e B recebe, portanto, um resistor externo de 4,7\,k$\Omega$ até 3,3\,V: quando o transistor está aberto, o ESP lê alto; quando conduz, ele drena a pequena corrente do resistor e o ESP lê baixo. A alimentação do corpo do encoder é separada da tensão dos sinais. A documentação do modelo recomenda alimentação externa de 12\,V, mas o conector da revisão de esquema analisada sugere outra linha; conferir etiqueta, folha de dados e montagem antes de ligar. \textbf{Nunca ligar o pull-up dos sinais à alimentação de 12 V.} O GND dessa fonte e o GND do Mestre precisam ser comuns para que alto e baixo tenham referência definida.

O uso de \emph{dois encoders}, um em cada lado, permite estimar a distância média mesmo quando as rodas não percorrem exatamente o mesmo caminho em uma curva. Também permite comparar os lados: diferença persistente pode indicar deslizamento, geometria incorreta, interrupção de um cabo ou falha de um sensor. Essa comparação é diagnóstica e não identifica sozinha a causa. Cabos de sinal devem ser afastados dos de potência; confirmar blindagem, retorno de GND, forma de onda e ausência de contagem com rodas paradas.

O periférico PCNT conta as quatro bordas dos canais A e B de cada lado. Assim, 600 P/R produzem 2\,400 contagens por volta do eixo do encoder. O filtro de pulsos curtos configurado é de 1\,000\,ns. O odômetro é distância acumulada sem sinal, não posição orientada; tampouco decide a marcha do veículo.

Para cada lado $j$, a distância estimada é
\begin{equation}
d_j=\frac{n_j}{4P_j\,i_j}\,\pi D_j ,
\qquad
v_j=\frac{\Delta d_j}{\Delta t}\,3{,}6 ,
\label{eq:encoder}
\end{equation}
em que $n_j$ é a contagem acumulada, $P_j$ os pulsos nominais por volta, $i_j$ as voltas do encoder por volta da roda, $D_j$ o diâmetro efetivo em metros e $\Delta t$ o intervalo em segundos. O fator 3,6 converte m/s em km/h. O firmware configura $D=0{,}300$\,m e $i=1{,}0$ para cada lado; são valores a medir na montagem. O serviço usa a média das distâncias esquerda e direita, rejeita saltos acima do limite diagnóstico de 200\,km/h e alerta sobre divergência persistente. Esse limite é filtro de plausibilidade, \textbf{não} velocidade certificada do carro.

O percurso da sessão reinicia no boot. O odômetro total é recuperado da NVS e persistido a cada 100\,m, compromisso entre desgaste de flash e perda potencial após corte de energia. Medir circunferência efetiva com o pneu carregado, confirmar relação de transmissão e comparar ambos os lados com uma referência independente antes de usar distância ou velocidade em análise quantitativa.

\section{MPU6050: aceleração e rotação}
O MPU6050 reúne um \emph{acelerômetro}, que responde à aceleração e à gravidade, e um \emph{giroscópio}, que mede a velocidade de rotação. Ele conversa com o Mestre por I\textsuperscript{2}C no endereço \texttt{0x68}; SDA vai ao GPIO 8 e SCL ao GPIO 9. A placa do sensor deve ficar rigidamente fixada ao chassi. O firmware foi parametrizado para X vertical positivo para cima e Y longitudinal positivo para frente. Se a montagem girar ou inverter o módulo, é preciso ajustar eixo e sinal em \path{firmware/mestre-s3/include/nucleo/configuracao_placa.h} e repetir a calibração.

O acelerômetro não distingue sozinho a aceleração do veículo da componente da gravidade produzida por inclinação da placa. O serviço amostra a 100\,Hz, usa o giroscópio para acompanhar mudanças rápidas de orientação e corrige lentamente a estimativa de gravidade quando o movimento permite. Subtrai então a gravidade projetada no eixo longitudinal, filtra ruído e usa dois limiares com confirmação temporal para classificar aceleração, desaceleração ou condição próxima do repouso sem alternar a cada vibração. O valor em \emph{g} é convertido nominalmente por $a_{\mathrm{m/s^2}}=9{,}80665\,a_g$.

Ao iniciar, o veículo deve permanecer parado: o firmware usa 100 amostras para estimar os desvios e rejeita calibração com giro, intensidade de gravidade implausível ou eixo longitudinal quase vertical. Uma calibração aceita não corrige uma fixação que depois se moveu. A velocidade principal vem dos encoders, pois integrar a aceleração inercial para obter velocidade acumularia erro de inclinação, ruído e viés. Telemetria e diagnóstico distinguem sensor presente, calibração concluída e dado válido; conferir esses estados antes de comparar aceleração com outro instrumento.

\section{Tela, sinalização e armazenamento}
A OLED SSD1306 de 0,96\,pol e 128$\times$64 pixels usa SDA/SCL do mesmo barramento I\textsuperscript{2}C que atende o MPU6050. O firmware procura os endereços \texttt{0x3C} e \texttt{0x3D}; isso permite encontrar variantes comuns sem alterar o programa. Conferir alimentação e nível dos pull-ups no módulo adquirido: o barramento precisa permanecer compatível com 3,3\,V, e resistores de módulos em paralelo alteram sua resistência equivalente. O capítulo~\ref{cap:eletronica} explica o dimensionamento. A tela apresenta etapas de partida, estado, medições disponíveis e progresso de operações; sua ausência é reportada, mas não impede as UARTs nem o registro. As páginas e os avisos são detalhados no capítulo~\ref{cap:diagnostico}.

Os quatro LEDs externos mostram estados resumidos sem exigir telefone ou terminal: verde para condição normal, amarelo para atenção, vermelho para falha crítica e azul para fases como partida, Wi-Fi e OTA. Cada um precisa de resistor em série dimensionado para a corrente e luminosidade do componente instalado; o esquema/fiação deve confirmar sua polaridade. O buzzer passivo recebe tons por PWM no GPIO 14. Como o GPIO não deve alimentar uma carga acima de sua capacidade, prever transistor e resistor de base conforme o buzzer escolhido, com retorno em GND; se a carga for indutiva, prever proteção contra o pulso de desligamento. A porcentagem de PWM configurada não equivale a um nível acústico medido.

O microSD usa SPI2, com CLK, MOSI, MISO e CS nos GPIOs 38, 39, 40 e 41. Ele armazena CSV, logs e cópia de ajustes; uma tarefa de gravação recebe amostras por fila para que o tempo de escrita do cartão não dite o ritmo da aquisição. A falta do cartão elimina esses arquivos persistentes enquanto durar a falha, embora o fluxo ao vivo possa continuar. Confirmar alimentação, nível lógico, desacoplamento e pinagem no módulo comprado; um módulo com regulador ou conversores de nível pode ter requisitos diferentes de um encaixe de cartão simples. O capítulo~\ref{cap:memorias_google} descreve arquivos e recuperação.

\section{UART e referência elétrica entre os três ESPs}
\begin{table}[htbp]\centering\small
\caption{Ligações bidirecionais entre os três ESP32-S3}\label{tab:uart-fios}
\begin{tabularx}{\textwidth}{@{}p{2.2cm}X X p{1.8cm}@{}}\toprule
Enlace & Mestre TX $\to$ satélite RX & Satélite TX $\to$ mestre RX & Referência\\\midrule
Equipe & GPIO 17 $\to$ GPIO 5 & GPIO 6 $\to$ GPIO 18 & GND comum\\
Visitantes & GPIO 15 $\to$ GPIO 5 & GPIO 6 $\to$ GPIO 16 & GND comum\\\bottomrule
\end{tabularx}\end{table}
Cada satélite tem um par TX/RX exclusivo do Mestre. TX transmite e deve chegar ao RX da outra placa; unir TX com TX impede a comunicação. Os três GNDs precisam compartilhar uma referência elétrica para que cada receptor interprete corretamente os níveis dos sinais. \textbf{Não é necessário interligar as linhas de 3,3\,V das três placas}: cada uma pode receber alimentação pelo seu ramo previsto, desde que os sinais e retornos respeitem os limites elétricos. Aterramento comum não significa usar o fio fino da UART como caminho de retorno da corrente de potência.

Ambos os enlaces usam 1\,000\,000 baud, oito bits de dados, sem paridade e um bit de parada (8N1). O Mestre envia telemetria e recebe confirmações; a Equipe também encaminha solicitações, e os satélites podem receber imagens OTA pelo mesmo enlace. Antes de interpretar erro de protocolo, conferir alimentação de cada placa, continuidade do GND, cruzamento TX/RX, nível alto de 3,3\,V e integridade do cabo. Se as placas estiverem fisicamente afastadas ou em ambiente ruidoso, verificar a forma de onda e o aterramento na montagem real.

\chapter{Eletrônica da placa e instrumentação elétrica}\label{cap:eletronica}
\section{Alcance do esquema analisado}
As folhas hierárquicas consultadas se identificam como revisão v4 de 18/10/2026. Elas representam o circuito projetado; o repositório não associa essa revisão a uma PCB comprovadamente montada no veículo. O símbolo do mestre usa a designação ESP32-S3-DevKitC, enquanto o firmware é configurado para o módulo N16R8. Antes de fabricar ou modificar a placa, conferir módulo, footprint, pinos físicos, esquema, lista de materiais e identificação impressa da PCB. Valores e ligações citados abaixo descrevem o esquema, salvo quando o texto identifica uma implementação no firmware ou um requisito de ensaio.

\section{Cadeia de alimentação: 48 V, 5 V e 3,3 V}
A folha de alimentação representa uma entrada rotulada 48 V, um módulo buck identificado como 100 V/5 V--2 A, saída de 5 V com C1 de 47\,$\mu$F e um estágio indicado como AMS1117 para produzir 3,3 V. O buck é um conversor chaveado abaixador: transfere energia por uma chave e um indutor para obter 5 V com perdas menores que uma queda linear direta desde o barramento do veículo. \textbf{Na distribuição projetada, a linha de 5 V alimenta os módulos ESP e outros consumidores de 5 V; a linha de 3,3 V fornece alimentação a determinados periféricos e os pull-ups dos sinais.} Os pinos GPIO dos ESPs continuam sendo sinais lógicos de 3,3 V, mesmo quando a placa de desenvolvimento recebe 5 V no pino de alimentação.

O rótulo de 100 V/5 V--2 A identifica o símbolo customizado, não uma medição da capacidade do módulo comprado. Confirmar fabricante, faixa de entrada, tensão máxima do conjunto de baterias carregado, corrente contínua e de pico, dissipação, proteção contra surtos e polaridade. Para o estágio linear, a potência térmica aproximada é $P_{\mathrm{LDO}}=(V_{\mathrm{in}}-V_{\mathrm{out}})I$; de 5 V para 3,3 V, cada 0,5 A requer dissipar cerca de 0,85 W. É preciso somar os consumos de todos os ramos, inclusive os picos de Wi-Fi e escrita no SD, e medir as duas linhas durante a partida. O fato de uma placa ESP aceitar 5 V na entrada não permite alimentar seus GPIOs com 5 V.

\begin{figure}[htbp]\centering
\begin{tikzpicture}[>=Latex,bloco/.style={draw=nexus,rounded corners,fill=pale,align=center,minimum height=10mm,font=\small}]
\node[bloco,text width=2.2cm] (bat) at (0,0) {Entrada rotulada\\48 V};
\node[bloco,text width=2.2cm] (buck) at (3.3,0) {Buck\\5 V};
\node[bloco,text width=3.0cm] (cinco) at (7.3,0) {Placas ESP e módulos\\alimentados em 5 V};
\node[bloco,text width=2.5cm] (ldo) at (3.3,-1.9) {Estágio AMS1117\\3,3 V};
\node[bloco,text width=3.0cm] (tres) at (7.3,-1.9) {Módulos de 3,3 V\\e pull-ups};
\draw[->,thick] (bat)--(buck);
\draw[->,thick] (buck)--(cinco);
\draw[->,thick] (buck)--(ldo);
\draw[->,thick] (ldo)--(tres);
\end{tikzpicture}
\caption{Distribuição de alimentação representada no esquema; confirmar cada saída na placa montada.}\label{fig:alimentacao}
\end{figure}

\begin{table}[htbp]\centering\small
\caption{Linhas de alimentação e cargas indicadas no esquema}\label{tab:alimentacao-modulos}
\begin{tabularx}{\textwidth}{@{}p{3.6cm}p{2.0cm}X@{}}\toprule
Conjunto & Linha indicada & O que conferir na montagem\\\midrule
Placas ESP mestre, Equipe e Visitantes & 5 V na placa & Entrada de 5 V e regulador interno da placa usada; GPIO e UART permanecem em 3,3 V.\\
Buzzer, circuito do relé e conector do módulo INA228 & 5 V & Corrente da bobina, módulo realmente instalado e compatibilidade dos sinais com 3,3 V.\\
MPU6050, OLED e microSD & 3,3 V & Tensões nos conectores, consumo e desacoplamento local.\\
SDA/SCL e canais A/B dos encoders & 3,3 V via resistores & Pull-ups são referências de nível lógico, não fontes de alimentação dos módulos.\\
Corpo dos encoders & Divergente & O conector do esquema indica 5 V; a documentação do modelo cita alimentação externa de 12 V. Conferir etiqueta e folha de dados do componente instalado antes de ligar.\\\bottomrule
\end{tabularx}\end{table}

O esquema apresentado não fecha, por si só, uma análise de proteção da entrada: fusível, proteção contra inversão, sobretensão transitória, desconexão de manutenção e segregação de cabos de tração precisam ser identificados na placa real e nos documentos do veículo. O GND comum é necessário para sinais UART e I\textsuperscript{2}C, mas correntes altas do motor não devem circular pelo retorno fino dos sensores. Planejar retornos e pontos de medição para evitar que queda de tensão em trilhas seja interpretada como sinal.

\section{Desacoplamento e transientes}
Na folha de alimentação, C1 de 47\,$\mu$F aparece entre 5 V e GND na saída do buck. Na folha de periféricos, há 100\,$\mu$F e 100\,nF junto da alimentação de 3,3 V do módulo microSD. O capacitor maior fornece reserva local para variações mais lentas de corrente; o cerâmico de 100\,nF fornece um caminho de baixa impedância para componentes rápidos. Eles se complementam: um valor grande isolado não substitui um capacitor pequeno colocado perto do pino, e o capacitor pequeno não sustenta uma demanda longa.

Posicionar desacoplamentos com trilhas curtas de alimentação e retorno; medir afundamento e ruído nos pinos durante Wi-Fi, escrita SD e acionamentos. Verificar tensão nominal, corrente de ripple, polaridade e ESR (resistência série equivalente) conforme o conversor e o regulador efetivamente usados. Os módulos ESP e sensores podem incluir capacitores próprios; sua presença e valor não devem ser inferidos apenas pelo símbolo do conector. As diretrizes de hardware do ESP32-S3 recomendam desacoplamento local também no domínio de 3,3 V e junto à memória.

\section{Barramento I\textsuperscript{2}C e resistores de pull-up}
I\textsuperscript{2}C usa SDA (dados) e SCL (clock) compartilhados. As saídas são do tipo \emph{dreno aberto}: um dispositivo consegue puxar a linha para zero, mas não a força ativamente para nível alto. Um resistor de \emph{pull-up} em cada linha fornece esse alto quando ninguém a puxa para baixo. No esquema v4, R5 e R6 são de 4,7\,k$\Omega$ para 3,3 V; o firmware configura 400\,kHz no mestre. MPU6050 e OLED compartilham o barramento; o conector J8 reserva SDA/SCL para o módulo associado ao INA228. Endereços diferentes permitem coexistência, mas o endereço efetivo do INA228 depende da montagem de A0/A1 no módulo e ainda deve ser registrado.

Módulos comerciais frequentemente já trazem resistores de pull-up. Todos os resistores da mesma linha ficam em paralelo:
\begin{equation}
\frac{1}{R_{\mathrm{eq}}}=\sum_{k=1}^{n}\frac{1}{R_k}.
\label{eq:pullup}
\end{equation}
Duas unidades de 4,7\,k$\Omega$ equivalem a 2,35\,k$\Omega$; quatro equivalem a 1,175\,k$\Omega$; cinco, a 940\,$\Omega$. Resistência excessivamente baixa exige que os dispositivos drenem corrente alta para produzir zero; resistência alta demais, junto à capacitância dos cabos e módulos, faz a borda subir lentamente. A janela admissível depende de $V_{\mathrm{DD}}$, corrente de afundamento, capacitância total e tempo de subida. O limite inferior aproximado é $R_{p,\min}=(V_{\mathrm{DD}}-V_{OL,\max})/I_{OL}$: se o dispositivo suportar 3\,mA com nível baixo máximo de 0,4\,V a 3,3\,V, resulta cerca de 967\,$\Omega$. Cinco resistores de 4,7\,k$\Omega$ em paralelo já ficariam abaixo desse exemplo. Para modo rápido a 400\,kHz, a nota de aplicação da TI dá $t_r \simeq 0{,}8473R_{\mathrm{eq}}C_b$ e limite de subida de 300\,ns. Se $C_b=100$\,pF, o limite superior calculado é cerca de 3,54\,k$\Omega$; 4,7\,k$\Omega$ sozinho pode ser lento, enquanto 2,35\,k$\Omega$ é um candidato. Em $C_b=200$\,pF, o limite superior cai para 1,77\,k$\Omega$. Esses exemplos não substituem medir o barramento real.

Antes de montar todos os módulos, identificar os pull-ups existentes em cada placa, medir a resistência equivalente de SDA e SCL com a alimentação desligada e decidir quais resistores de módulo devem ser removidos ou desconectados. \textbf{Manter ao menos um par de pull-ups a 3,3 V}, mas remover os pares excedentes quando a associação em paralelo sair da janela calculada. Se um módulo puxar SDA/SCL para 5 V, remover seus pull-ups para 5 V ou instalar adaptação de nível apropriada antes de ligar ao ESP32-S3. Confirmar o resultado no osciloscópio, com todos os módulos e cabos, durante transferências a 400\,kHz. Resistores dos encoders são outro circuito e não entram nessa conta.

\section{INA228 e shunt de 1 m$\Omega$}
O conector J8 da folha de periféricos é denominado INA228 e entradas do motorista; ele traz alimentação indicada como 5 V, GND, SDA/SCL e outros sinais. O firmware mestre contém um driver para o INA228 no endereço I\textsuperscript{2}C \texttt{0x40} e publica grandezas elétricas quando a leitura é válida. O código configura shunt nominal de 1\,m$\Omega$ e corrente máxima de 75\,A. A posição física do shunt, o módulo específico e um ensaio de calibração ainda precisam ser confirmados na montagem. Os parâmetros a seguir descrevem a intenção de projeto e a configuração do software, não uma leitura experimental já comprovada.

\begin{figure}[htbp]\centering
\begin{tikzpicture}[>=Latex,bloco/.style={draw=nexus,rounded corners,fill=pale,align=center,minimum height=9mm,font=\small}]
\node[bloco,text width=2.6cm] (fonte) at (0,0) {Banco de baterias\\positivo};
\node[bloco,text width=2.0cm] (shunt) at (3.5,0) {Shunt\\1 m$\Omega$};
\node[bloco,text width=2.5cm] (carga) at (7.0,0) {Controlador\\do motor};
\node[bloco,text width=2.4cm] (ina) at (3.5,-2.2) {INA228\\interface I\textsuperscript{2}C};
\draw[->,very thick] (fonte)--(shunt);
\draw[->,very thick] (shunt)--(carga);
\draw[->,accent,thick] (shunt.south west)--node[left,font=\scriptsize]{IN+}(ina.north west);
\draw[->,accent,thick] (shunt.south east)--node[right,font=\scriptsize]{IN$-$}(ina.north east);
\node[bloco,text width=2.5cm] (vbus) at (7.0,-2.2) {Ponto VBUS\\após o shunt};
\draw[->,accent,dashed] (carga.south)--(vbus.north);
\draw[->,accent,dashed] (vbus.west)--(ina.east);
\end{tikzpicture}
\caption{Exemplo de medição no lado positivo com VBUS após o shunt. As tomadas Kelvin e o ponto VBUS devem ser definidos no projeto elétrico; o desenho não comprova o chicote instalado.}\label{fig:shunt-kelvin}
\end{figure}

O INA228 recebe duas medições distintas. \textbf{VSHUNT} é a diferença entre IN+ e IN$-$, tomada diretamente nos extremos do elemento resistivo por fios Kelvin; dividida pela resistência conhecida, fornece a corrente e seu sentido. \textbf{VBUS} é a tensão entre o pino VBUS e GND no ponto positivo escolhido do barramento do banco de baterias. Não é a queda no shunt: ela indica a tensão do barramento que alimenta o controlador. Se o pino for conectado antes ou depois do shunt, a leitura difere aproximadamente pela própria queda no shunt; a posição precisa constar do esquema final. A potência elétrica pode ser calculada a partir de corrente e tensão de barramento coerentes no tempo, $P\simeq V_{\mathrm{BUS}}I$, com convenção de sinal definida para consumo e regeneração. O driver no Mestre já lê os registradores de VSHUNT, VBUS, corrente, potência, temperatura do chip, energia e carga acumuladas. A validade exige resposta do dispositivo e verificações de faixa; uma leitura disponível ainda requer calibração experimental da montagem.

Com $R_{\mathrm{shunt}}=1\,\mathrm{m}\Omega=0{,}001\,\Omega$:
\begin{equation}
V_{\mathrm{shunt}}=IR_{\mathrm{shunt}},\qquad
1\,\mathrm{A}\longleftrightarrow 1\,\mathrm{mV},\qquad
75\,\mathrm{A}\longleftrightarrow 75\,\mathrm{mV}.
\label{eq:shunt}
\end{equation}
O calor no elemento resistivo é $P_{\mathrm{shunt}}=I^2R$: em 75 A, 5,625 W. Um shunt anunciado para 75 A precisa suportar essa dissipação no ambiente e na fixação reais, com margem térmica, seção dos terminais e tempo de sobrecarga adequados. A resistência muda com temperatura e com a tolerância de fabricação. As duas linhas de medição devem chegar a pontos Kelvin próximos ao elemento resistivo, separados do trajeto de corrente de potência, para não somar queda de tensão dos parafusos e cabos.

Segundo a folha de dados da TI, o INA228 oferece faixas diferenciais de $\pm163{,}84$\,mV e $\pm40{,}96$\,mV em VSHUNT. Para 75 A com shunt de 1\,m$\Omega$, a faixa de $\pm40{,}96$\,mV saturaria antes da corrente nominal; deve-se selecionar a faixa ampla (\texttt{ADCRANGE=0}) se a meta for medir até 75 A. Esse nome identifica a seleção do conversor interno do INA228. VBUS aceita de 0 a 85 V segundo a folha de dados do circuito integrado; a tensão máxima do banco de baterias, inclusive sob carga completa e transitórios, precisa respeitar o limite do circuito e do módulo comprado. A alimentação lógica VS e os pull-ups SDA/SCL são assuntos distintos: o conector J8 indica 5 V, mas isso não autoriza aplicar 5 V aos GPIOs do mestre. Verificar regulador, tradutor de nível e resistores da placa INA228 adquirida.

Para converter o registrador CURRENT em ampères, o firmware configura a calibração do INA228 com a resistência nominal do shunt. A folha de dados define, para \texttt{ADCRANGE=0},
\begin{equation}
\mathrm{CURRENT\_LSB}\geq \frac{I_{\max}}{2^{19}},
\qquad
\mathrm{SHUNT\_CAL}=13107{,}2\cdot10^6\,
       \mathrm{CURRENT\_LSB}\,R_{\mathrm{shunt}} .
\label{eq:ina-cal}
\end{equation}
No código, \texttt{CURRENT\_LSB} é calculado como $75/2^{19}$\,A por bit, aproximadamente 143\,$\mu$A, e \texttt{SHUNT\_CAL} resulta em cerca de 1875 para 1\,m$\Omega$. O Mestre confere identificação do chip e leitura de retorno da calibração, usa a faixa ampla de VSHUNT e rejeita erro matemático, conversão ainda não pronta, temperatura fora da faixa e corrente acima do limite admitido. Após três erros seguidos, suspende a leitura e tenta reinicializar o módulo após cinco segundos. A telemetria marca grandezas indisponíveis quando não há leitura válida; o simulador não deve ser confundido com o sensor. Definir a polaridade IN+/IN$-$ para que consumo e regeneração tenham sinais inequívocos. Comparar leituras em vários pontos com instrumento de referência, inclusive perto de zero e perto do máximo. Um alerta de sobrecorrente por software não substitui proteção elétrica independente.

\section{Relé habilitador, botão de emergência e diodo de roda livre}
A função prevista para K1 é interromper a \emph{linha de habilitação} do driver do motor elétrico, e não conduzir a corrente de tração. A intenção de integração é pôr um contato normalmente aberto do relé em série com um botão de emergência normalmente fechado, no circuito de habilitação exigido pelo driver. Com ambos fechados, o driver pode ser habilitado; ao pressionar o botão, abrir o contato ou perder a alimentação da bobina, a linha é interrompida. Assim, o condutor dispõe de uma ação local e o sistema poderia pedir uma desabilitação remota. O sentido elétrico e os níveis da entrada ENABLE devem ser confirmados na documentação do driver utilizado; o diagrama abaixo representa a função pretendida, não um chicote já montado.

\begin{figure}[htbp]\centering
\begin{tikzpicture}[>=Latex,bloco/.style={draw=nexus,rounded corners,fill=pale,align=center,minimum height=10mm,font=\small}]
\node[bloco,text width=2.3cm] (sinal) at (0,0) {Fonte do sinal\\de habilitação};
\node[bloco,text width=2.5cm] (botao) at (3.1,0) {Botão de emergência\\NF};
\node[bloco,text width=2.3cm] (contato) at (6.2,0) {Contato K1\\NA};
\node[bloco,text width=2.4cm] (driver) at (9.3,0) {Entrada ENABLE\\do driver};
\draw[->,thick] (sinal)--(botao);
\draw[->,thick] (botao)--(contato);
\draw[->,thick] (contato)--(driver);
\node[bloco,text width=2.3cm] (gpio) at (0,-2.0) {Saída do ESP\\futura};
\node[bloco,text width=2.5cm] (base) at (3.1,-2.0) {Resistor de base\\e estado seguro};
\node[bloco,text width=2.3cm] (trans) at (6.2,-2.0) {Transistor\\Q4};
\node[bloco,text width=2.4cm] (bobina) at (9.3,-2.0) {Bobina K1\\D1 em paralelo};
\draw[->,thick,accent] (gpio)--(base);
\draw[->,thick,accent] (base)--(trans);
\draw[->,thick,accent] (trans)--(bobina);
\draw[densely dashed,->] (bobina.north)--node[right,font=\scriptsize]{aciona}(contato.south east);
\end{tikzpicture}
\caption{Duas partes da função pretendida: caminho de habilitação em série, acima, e acionamento da bobina, abaixo. NF: normalmente fechado; NA: normalmente aberto.}\label{fig:rele-habilitador}
\end{figure}

A folha do esquema contém K1 (símbolo SANYOU SRD), Q4 (BC337), R1 de 1\,k$\Omega$, R2 de 10\,k$\Omega$ e D1 (1N4007). \textbf{A função de desabilitação remota ainda não está implementada no firmware consultado}: o painel e o mestre guardam um pedido de parada na RAM, mas não acionam Q4, K1 ou uma entrada do driver. Para transformar o pedido em atuação, serão necessários mapeamento do GPIO, circuito que desenergize K1 em falha ou reinício, política para perda de comunicação, retorno do estado real do contato/driver e ensaio de resposta. O botão físico deve funcionar independentemente do navegador, Wi-Fi e UART.

A bobina é indutiva. Ao interromper sua corrente, o campo magnético tenta mantê-la e pode produzir sobretensão no transistor. O diodo de roda livre D1 em paralelo oferece um caminho temporário para essa corrente: cátodo no lado positivo da bobina e ânodo no lado comutado pelo transistor, ficando reversamente polarizado durante o acionamento normal. Ele reduz o pico de tensão, mas pode aumentar o tempo de soltura do relé; medir esse tempo no circuito final. Q4 evita exigir do GPIO a corrente da bobina; R1 limita a corrente de base, e o estado de desligamento da base precisa ser definido durante reset e partida.

\begin{nota}{Conferência obrigatória da folha v4}
No desenho exportado da folha, R2 de 10\,k$\Omega$ aparece em série com a alimentação da bobina e Q4 como derivação desse ramo; essa ligação não corresponde diretamente ao acionamento por transistor em série mostrado funcionalmente acima. A queda de tensão em R2 pode impedir a corrente nominal da bobina. É necessário revisar a ligação no esquema, conferir bobina, contato escolhido, orientação de D1 e estado sem alimentação, executar a checagem elétrica do KiCad e testar com fonte limitada em bancada antes de conectar o driver. A posição do botão de emergência em série também deve constar do esquema e do chicote finais.
\end{nota}

Outros transistores da folha representam sinais de modo e reversão e o MCP4725 aparece como DAC (conversor digital-analógico) para a referência de acelerador. Esses caminhos ainda não são acionados pelo controle preparatório implementado. Antes de integração, mapear tensão exigida pelo driver, estado seguro de cada fio, isolação, tempo de resposta, limites físicos e retorno medido; testar separadamente de toda a tração.

\paragraph{Fontes técnicas.}
As faixas e equações do INA228 vêm da \href{https://www.ti.com/lit/ds/symlink/ina228.pdf}{folha de dados da Texas Instruments}. O cálculo da janela de pull-ups segue \href{https://www.ti.com/lit/an/slva689/slva689.pdf}{I\textsuperscript{2}C Bus Pull-Up Resistor Calculation}, da TI. O desacoplamento do ESP deve ser confrontado com o \href{https://docs.espressif.com/projects/esp-hardware-design-guidelines/en/latest/esp32s3/schematic-checklist.html}{checklist de esquema do ESP32-S3}, da Espressif, e com a folha de dados de cada módulo comprado.

\chapter{Firmware e processamento}\label{cap:firmware}
\section{Projetos e responsabilidades do código}
O NEXUS tem três projetos PlatformIO/ESP-IDF independentes: mestre, Equipe e Visitantes. Cada ESP executa sua própria aplicação; o código de estruturas e protocolos compartilhados fica em \path{firmware/compartilhado}. O ambiente do mestre é \texttt{mestre-s3-n16r8} e fixa \texttt{espressif32@6.13.0}; os satélites usam \texttt{esp32-s3-devkitm-1}, sem fixar essa versão nos respectivos \texttt{platformio.ini}. Para reproduzir um ensaio, registrar também a versão efetivamente instalada das dependências dos satélites.

\begin{table}[htbp]\centering\small
\caption{Quem define e executa cada responsabilidade no mestre}\label{tab:camadas}
\begin{tabularx}{\textwidth}{@{}p{3.1cm}X@{}}\toprule
Parte do projeto & Responsabilidade e limite\\\midrule
\texttt{include/nucleo} & Define pinos, constantes, limites e política das tarefas. É a configuração usada pelos módulos, não um driver nem registro de teste da placa.\\
\texttt{src/gerenciadores} & Controla recursos do ESP-IDF e oferece operações de UART, I\textsuperscript{2}C, PCNT, SPI/cartão, Wi-Fi, OTA local e sinalização. Devolve erros para que a aplicação decida o que fazer.\\
\texttt{src/dispositivos} & Interpreta os sinais e registradores dos componentes conectados, como MPU6050, encoders e SSD1306. Uma leitura de dispositivo não decide sozinha o estado de saúde do sistema.\\
\texttt{src/servicos} & Combina dispositivos e gerenciadores para formar aceleração, odômetro, telemetria, histórico, comandos, atualização, supervisão e diagnóstico; cada serviço mantém seu estado e suas regras.\\
\texttt{src/main.cpp} & Compõe os serviços na partida, classifica falha essencial ou degradável, cria o ciclo de telemetria e coordena a observação da imagem OTA.\\
\texttt{firmware/compartilhado} & Define tipos de dados e contratos binários que mestre e satélites precisam interpretar do mesmo modo.\\\bottomrule
\end{tabularx}\end{table}

Nos satélites, \texttt{main.c} prepara os serviços locais; o receptor UART valida e distribui quadros; o servidor HTTP entrega páginas e respostas JSON aos navegadores. A Equipe encaminha consultas e ações ao mestre e dispõe de diagnóstico mais amplo. Visitantes oferece leitura pública. O servidor HTTP constitui uma API de rede nas rotas \texttt{/api/...}; uma função interna do firmware ou um quadro UART não é, por si só, uma API HTTP.

\section{Sequência de partida do mestre}
Ao ligar, o \emph{bootloader} escolhe uma partição de aplicação e transfere o controle a \texttt{app\_main}. O aplicativo registra a nova partida para a guarda de rollback e confere se a flash detectada corresponde ao alvo N16R8. Flash incompatível impede a execução. A PSRAM é verificada separadamente: ausência ou tamanho inesperado gera aviso e caminhos que a usam podem tentar RAM interna, com menos margem para buffers e histórico.

A sequência seguinte é intencional. Primeiro iniciam LEDs/buzzer e OLED, para informar as etapas já durante a partida. Em seguida abre-se a NVS, que contém odômetro e perfis de rede; se ela exigir manutenção ou estiver inacessível, o firmware não a apaga automaticamente para evitar perder dados. Depois inicia a supervisão, que acompanha o progresso das tarefas, e o núcleo da aplicação: serviço de telemetria, enlaces exclusivos com Equipe e Visitantes, simulador, MPU6050, encoders e odômetro. Falha de um dispositivo de medição pode deixar apenas aquela grandeza indisponível. Se os encoders não iniciarem, o odômetro não é iniciado; outros serviços podem continuar.

O programa tenta montar o microSD e iniciar o histórico, cria a tarefa de telemetria e a registra na supervisão. Sem a tarefa ou sem a supervisão essencial, reinicia para evitar um sistema parcialmente ativo. Com ela pronta, inicia Wi-Fi externo, OTA, menu serial, relatórios e diagnóstico. Falha de cartão, tela, sinalização ou serviço opcional é comunicada e permite continuar o que estiver funcional. Wi-Fi e OTA têm uma condição adicional: se uma imagem recém-atualizada ainda estiver em autoexame, falha em qualquer um desses serviços impede confirmá-la, pois eles integram o caminho de recuperação remota. O aplicativo então solicita retorno à versão anterior. O procedimento completo de validação e rollback está no Capítulo~\ref{cap:ota}.

O ciclo de telemetria acorda normalmente a cada 250\,ms. A tarefa do MPU6050 amostra a 100\,Hz; filas deixam transmissão UART e gravação no SD avançarem em tarefas próprias. A política fixa aquisição e enlaces principais no núcleo 1, enquanto OTA, CSV e upload ficam no núcleo 0; algumas tarefas não têm afinidade. São escolhas de escalonamento, não garantia de prazo absoluto. O diagnóstico mede atrasos, CPU, memória e margem de pilha durante cargas simultâneas.

\section{Formação da telemetria}
O ciclo do mestre monta um retrato \texttt{DadosTelemetriaVeiculo}. Na fonte real, obtém do serviço inercial as componentes, o estado de calibração e a aceleração longitudinal; atualiza odômetro e velocidade com os dois encoders; incorpora as grandezas elétricas do INA228 quando sua leitura é válida. Na fonte simulada, um gerador preenche o mesmo tipo de retrato e marca a origem como simulada. A troca de fonte limpa o retrato anterior para evitar misturar valores das duas origens. O padrão após reiniciar vem da configuração compilada; a Equipe pode alterá-lo temporariamente durante a sessão.

Antes do envio, o serviço rejeita valores não finitos e campos fora dos limites de plausibilidade definidos no código. Converte números do retrato em inteiros escalonados para manter um pacote binário de tamanho fixo. Exemplos do contrato são centésimos de km/h para velocidade, milésimos de m/s\textsuperscript{2} para aceleração e milivolts ou centésimos de ampère para campos elétricos quando estes tiverem fonte válida. O formato também reserva pedais, estado de bateria, localização, distâncias e autonomia. A presença desses campos no pacote não significa que os respectivos sensores estejam instalados ou integrados; sem origem real, a interface deve indicar indisponibilidade. A marcha solicitada pelo controle preparatório é um estado separado da marcha publicada.

O pacote recebe sequência crescente e tempo relativo ao início do mestre. Uma cópia segue por cada UART dedicada, outra alimenta relatórios e o registro histórico seleciona amostras a 1\,Hz. O bit mais alto do campo de marcha identifica telemetria simulada para os receptores. A ordem, a escala e a verificação dos bytes são um \emph{protocolo binário} compartilhado, não uma chamada HTTP. A seção seguinte detalha como o receptor distingue mensagens e confirma a chegada.

\section{Filtragem, odômetro e calibração}
O MPU6050 entrega aceleração e rotação. Durante a partida, o serviço só acumula amostras para calibração quando a intensidade da aceleração é compatível com a gravidade e o giroscópio indica repouso. Essa fase estima gravidade e desvio do giroscópio na orientação instalada. Se o eixo escolhido para representar o sentido longitudinal estiver quase alinhado à gravidade, a calibração não deve mascarar o erro de montagem. Em operação, o giroscópio acompanha mudanças de orientação, o vetor estimado de gravidade é retirado da aceleração medida e o componente longitudinal passa por filtro passa-baixas de 3\,Hz. Uma histerese com tempo de confirmação evita que pequenas oscilações troquem continuamente a indicação de aceleração/frenagem. O valor filtrado é uma estimativa de movimento longitudinal, sujeita à fixação, orientação, vibração e calibração em repouso.

Cada encoder fornece pulsos que o PCNT conta por hardware. O serviço lê a diferença desde a atualização anterior, converte contagens em giro de roda usando resolução, diâmetro e relação mecânica configurados e calcula distância e velocidade de cada lado. A velocidade e o percurso divulgados usam as duas rodas, com tratamento de saltos e divergência persistente; um salto que implicaria velocidade superior ao limite diagnóstico de 200\,km/h é rejeitado. O código atual soma bordas sem inferir o sentido de rotação: marcha à ré também aumenta distância acumulada. O percurso de sessão recomeça após o boot, enquanto o odômetro total é recuperado da NVS e persistido a cada 100\,m. Uma queda de energia pode perder a distância ainda não persistida. Diâmetro efetivo e relação encoder/roda registrados no firmware são hipóteses até a medição no veículo; o capítulo~\ref{cap:hardware} mostra como conferir esses parâmetros.

\section{Enlaces, confirmações e interfaces}
Mestre--Equipe e mestre--Visitantes são enlaces UART exclusivos, bidirecionais, em lógica de 3,3\,V, 1\,000\,000 baud e formato 8N1. Ambos transportam várias famílias de quadros. Um analisador comum reconhece o marcador, confere versão, tamanho e CRC antes de entregar o conteúdo ao serviço responsável. Bytes de números com mais de uma posição seguem a ordem \emph{little-endian}. Esse enquadramento evita que um bloco OTA seja interpretado como telemetria ou que uma confirmação de recebimento pareça resposta de comando.

\begin{table}[htbp]\centering\small
\caption{Famílias do protocolo UART entre mestre e satélites}\label{tab:familias}
\begin{tabularx}{\textwidth}{@{}p{2.6cm}p{2.4cm}X@{}}\toprule
Família e início & Sentido usual & O que o receptor verifica e confirma\\\midrule
Telemetria v4, \texttt{AA 55} & Mestre para cada satélite & Pacote fixo de 105 bytes, com carga de 98 bytes e CRC-16/CCITT-FALSE; após validar conteúdo, atualiza o retrato e sua sequência.\\
ACK v4, \texttt{AA 56} & Cada satélite para mestre & Confirmação de 20 bytes com sequência e tempo do mestre. Permite acompanhar recebimento e tempo de ida e volta (RTT).\\
Comando v1, \texttt{C7 7C} & Solicitação e resposta & Identificador, origem, destino, resultado e CRC-32. A resposta memorizada para ID repetido evita executar uma mesma ação duas vezes.\\
OTA UART v1, \texttt{A5 5A} & Mestre e satélite alvo & Sessão, sequência, deslocamento, CRC-32 por quadro e SHA-256 da imagem completa; cada bloco é confirmado antes de avançar.\\\bottomrule
\end{tabularx}\end{table}

O satélite distingue uma sequência nova, perdida, repetida ou fora de ordem. Uma repetição pode receber novo ACK sem criar uma segunda amostra. O tempo relativo do mestre ajuda a reconhecer um reinício e retomar a sequência. CRC detecta alteração acidental, mas não autentica o emissor. O ACK comprova recepção de telemetria pelo satélite; o resultado de um comando, de um upload ou de uma instalação OTA é acompanhado por seus estados próprios. Os campos de versão e capacidades permitem que firmwares de datas diferentes rejeitem operações não suportadas em vez de interpretar cargas incompatíveis.

O navegador não participa do barramento UART. Ele chama rotas HTTP do ESP da Equipe ou de Visitantes, por exemplo \texttt{GET /api/telemetria}. Essas rotas formam a API HTTP dos painéis. A Equipe transforma algumas consultas e ações HTTP em comandos do protocolo UART e devolve um resultado ao navegador. A arquitetura, o acesso local e as páginas são detalhados nos capítulos~\ref{cap:comandos} e~\ref{cap:wifi-paineis}.

\section{Histórico e armazenamento}
O serviço de registro escolhe uma amostra por segundo do retrato que circula a 4\,Hz, usa fila de 24 posições e deixa uma tarefa proprietária da escrita no cartão. O arquivo ativo é sincronizado aproximadamente a cada 10\,s e fechado após uma hora ou 8\,MiB. Uma fila cheia aparece em contadores de perda; ela não atrasa diretamente a aquisição. Se o cartão não montar, o ciclo UART continua e a aplicação tenta novamente após 30\,s.

Logs estruturados usam uma área distinta do CSV. O serviço mantém uma janela de saúde em PSRAM, ou menor em RAM interna, para registrar o contexto anterior a uma falha quando o cartão estiver disponível. O tempo civil só passa a acompanhar os registros depois de uma referência UTC obtida de resposta HTTPS válida; antes disso há tempo relativo ao boot. NVS guarda odômetro e preferências. O arquivo \path{/sdcard/configuracao_atual.txt} salva uma cópia seletiva com sequência e checksum aproximadamente a cada minuto, sem senhas. O capítulo~\ref{cap:memorias_google} descreve recuperação de arquivos, política de backup e confirmação dos envios ao Google.

\section{Diagnóstico, sinalização e recuperação}
O diagnóstico reúne resultados de dispositivos, erros I\textsuperscript{2}C e UART, idade das confirmações dos satélites, progresso das tarefas, RAM interna, PSRAM, margem de pilha, Wi-Fi e falhas OTA. Atualiza o resumo de saúde a cada segundo e conserva contadores para distinguir uma causa presente de um erro que já cessou. A supervisão usa pulsos de progresso e prazos para detectar tarefas atrasadas; recuperações potencialmente bloqueantes executam fora da tarefa de alta prioridade. Cada enlace UART tem recuperação própria, sem reiniciar necessariamente o outro satélite.

\begin{table}[htbp]\centering\small
\caption{Estados publicados pela saúde do mestre e indicação ao operador}\label{tab:estados-saude-firmware}
\begin{tabularx}{\textwidth}{@{}p{2.4cm}X p{4.4cm}@{}}\toprule
Estado & Significado no resumo de diagnóstico & Sinalização local típica\\\midrule
Saudável & Nenhuma causa ativa no ciclo de avaliação. & LED verde contínuo; OLED nas páginas normais.\\
Atenção & Uma causa ativa não classificada como crítica, por exemplo sensor, barramento, enlace, supervisão, Wi-Fi fraco ou falha OTA recente. As funções que não dependem da causa podem continuar. & LED amarelo periódico; OLED mostra a causa e alterna com as páginas usuais.\\
Falha & Memória interna/PSRAM em nível crítico ou margem de pilha crítica; também pode haver falha de partida tratada pela aplicação antes de existir um resumo normal. & LED vermelho predominante com breve pausa; tela e serial exibem detalhe quando disponíveis.\\
Recuperando & A causa ativa desapareceu, mas o sistema conserva por cerca de 2\,s uma indicação de transição. & Indicação de aviso amarelo durante a transição, antes de voltar ao verde.\\\bottomrule
\end{tabularx}\end{table}

\emph{Operação degradada} descreve uma função indisponível enquanto outras continuam: por exemplo, sem microSD há telemetria ao vivo mas não há CSV persistente; sem OLED restam serial, LEDs e painéis; sem um satélite o outro pode seguir recebendo quadros. Esse termo não equivale automaticamente ao estado \emph{Falha} da tabela, que tem regras próprias de classificação. Já uma falha essencial de NVS, inicialização do núcleo, criação ou supervisão da tarefa de telemetria leva a rejeição da imagem OTA pendente e reinício para evitar execução incompleta.

Os quatro LEDs, o buzzer passivo e a OLED apresentam níveis diferentes de detalhe. Azul identifica partida, sessão Wi-Fi ou fase OTA; padrões de verde, amarelo e vermelho mostram saúde ou atualização. Sons curtos marcam eventos e transições, com prioridade para falhas; não são um alarme contínuo para cada tentativa Wi-Fi. A OLED apresenta progresso, páginas de telemetria e comunicação, mensagens e causas resumidas. O terminal serial e a API HTTP de diagnóstico do painel Equipe oferecem contadores e últimos erros, úteis para determinar o componente responsável. Um sensor não medido ou um campo reservado deve aparecer como indisponível na interface, não como diagnóstico inferido de um valor zero. O capítulo~\ref{cap:diagnostico} detalha condições, limiares e apresentação de cada causa.

Os pedidos de marcha, limites e parada disponíveis na página Controle pertencem a um estado preparatório em RAM. Eles não acionam GPIO, relé, motor ou freio no firmware examinado. Regras, limites e implicações de segurança estão concentrados no capítulo~\ref{cap:comandos}, para evitar duas descrições divergentes do mesmo controle.

\chapter{Comandos, consultas e controle preparatório}\label{cap:comandos}
\section{Canal de comandos e caminho da requisição}
A telemetria é uma publicação periódica do mestre. Operações como consultar a versão, trocar a fonte de dados, iniciar sincronização e autorizar OTA precisam de requisição e resposta: o operador deve saber se o pedido chegou, foi aceito ou foi rejeitado. Por isso, o enlace UART contém uma segunda família de quadros, além da telemetria. Essa separação evita interpretar um ACK de amostra como resultado de uma ação e permite prazos diferentes para uma consulta rápida e uma transferência demorada. O satélite da equipe é a ponte HTTP--UART; o visitante não recebe uma interface de operação.

\begin{figure}[htbp]\centering
\begin{tikzpicture}[node distance=20mm, >={Latex},
  bloco/.style={draw=nexus,rounded corners,fill=pale,align=center,
    minimum height=13mm,text width=3.3cm,font=\small}]
\node[bloco] (browser) {Navegador da Equipe};
\node[bloco,right=of browser] (equipe) {ESP Equipe\\servidor HTTP};
\node[bloco,right=of equipe] (master) {ESP Mestre\\serviço solicitado};
\draw[->,thick] ([yshift=3mm]browser.east)--node[above,font=\scriptsize]{requisição HTTP}([yshift=3mm]equipe.west);
\draw[->,thick] ([yshift=3mm]equipe.east)--node[above,font=\scriptsize]{comando UART1}([yshift=3mm]master.west);
\draw[->,thick,accent] ([yshift=-3mm]master.west)--node[below,font=\scriptsize]{resposta UART}([yshift=-3mm]equipe.east);
\draw[->,thick,accent] ([yshift=-3mm]equipe.west)--node[below,font=\scriptsize]{resposta HTTP}([yshift=-3mm]browser.east);
\end{tikzpicture}
\caption{Encaminhamento de uma requisição do painel técnico ao mestre. Operações demoradas têm estado consultado depois da resposta inicial.}\label{fig:comando}
\end{figure}

\subsection{Formato e verificações no UART}
O protocolo de comandos v1 começa por \texttt{C7 7C}. O cabeçalho compactado inclui versão, tipo solicitação/resposta, origem, destino, código de comando, identificador da solicitação, tamanho, resultado e campo reservado. A carga pode ter até 128 bytes e o quadro termina com CRC-32. O receptor verifica início, versão, sentido, tamanho e CRC antes de interpretar campos; números multibyte são codificados em \emph{little-endian}. O resultado explícito pode ser OK, não suportado, carga inválida, ocupado, negado, erro interno ou timeout. Uma resposta antiga ou de outro ID não deve ser usada como confirmação do pedido atual.

O mesmo ID pode aparecer novamente após perda da resposta. O mestre conserva respostas recentes e devolve a resposta memorizada sem repetir a ação com efeito colateral. Esse comportamento é importante para sincronização, mudanças de rede e OTA. As máscaras de capacidades anunciadas na resposta de versão e em consulta específica permitem que o painel oculte ou desabilite ações que um firmware mais antigo não entende. O mestre usa a capacidade anunciada pelos satélites também para escolher o tamanho compatível de bloco OTA. Um campo de capacidade é parte de compatibilidade de protocolo, não um indicador físico do veículo.

\begin{table}[htbp]\centering\small
\caption{Grupos de códigos do protocolo de comandos v1}\label{tab:codigos-comandos}
\begin{tabularx}{\textwidth}{@{}l l X@{}}\toprule
Códigos & Tipo & Finalidade\\\midrule
1--8 & Consultas básicas & Ping, versão, capacidades, saúde, telemetria, UART, OTA e Wi-Fi.\\
20--25 & OTA & Verificação, autorizações por alvo, cancelamento e atualização de todos.\\
30--33 & Mestre & Resumo no terminal, ação estruturada, configuração e rede salva.\\
34--40 & Diagnóstico e satélites & Ping/versão de satélite, diagnóstico, sincronização, simulador, clientes e gestão de AP.\\
41 & Controle preparatório & Consultar estado solicitado de marcha, limites e parada.\\\bottomrule
\end{tabularx}\end{table}

O código 31 contém uma ação tipada, com campos para número, índice, intervalo ou texto, em vez de criar um comando UART para cada botão do navegador. As ações implementadas incluem pausar relatórios, zerar percurso, sincronizar Google, ajustar intervalos, preparar/salvar/conectar redes, alterar política Wi-Fi, selecionar simulador e alterar pedidos preparatórios. Consultas de configuração e perfis salvos são separadas para não produzir efeito colateral.

\section{HTTP da Equipe e acesso local}
O satélite Equipe disponibiliza uma API HTTP local. O navegador usa \texttt{GET} em \path{/api/mestre?consulta=NOME} para solicitar dados do mestre e \texttt{POST} na mesma rota para enviar uma ação. O satélite traduz essas chamadas em comandos do protocolo UART; o mestre não expõe essa API por Wi-Fi. As rotas \path{/api/telemetria}, \path{/api/diagnostico} e \path{/api/wifi/clientes} consultam dados do próprio satélite. Não há chave administrativa HTTP adicional: quem se associa ao AP Equipe pode enviar ações locais ou dirigidas ao Mestre. O servidor reserva até duas consultas simultâneas ao Mestre e uma ação; chamadas adicionais podem receber HTTP 503. Consulta desconhecida ou método errado é rejeitado.

Algumas alterações de Wi-Fi exigem uma senha de permissão verificada pelo \emph{Mestre}. Ela vale somente para essas mudanças protegidas e não é autorização geral para os outros comandos. A senha do AP Equipe deve ser distribuída apenas a operadores porque o HTTP local não usa TLS; o AP Visitantes não oferece comandos. Não incluir a senha de permissão nem a senha do AP técnico em capturas, QR público ou exportação de telemetria.

\subsection{Resposta imediata versus término da operação}
Um \texttt{ok:true} ou resposta UART OK pode indicar que o mestre aceitou um pedido; não significa que um CSV chegou ao Google, que a imagem OTA reiniciou e foi validada, ou que uma alteração de rede completou. A interface consulta o estado posteriormente e apresenta progresso e erro final. Esse desenho evita um falso sucesso quando a operação continua em tarefa própria. Sem resposta, o valor anterior mostrado no navegador não confirma uma nova solicitação. Se o prazo esgotar, consultar diagnóstico, ID, tentativas e estado antes de repetir.

\section{Pedidos de marcha, limites e parada}
A página Controle mostra leituras recentes e os pedidos armazenados. Marcha solicitada admite neutro, frente ou ré; velocidade máxima solicitada usa centésimos de km/h (0--3500, equivalentes a 0--35\,km/h); aceleração máxima usa milésimos de m/s\textsuperscript{2} (0--5000); parada admite solicitar ou liberar. Na partida, a parada está solicitada, a marcha é neutra e ambos os limites são zero. Uma marcha diferente de neutro exige parada liberada e limites positivos. Zerar um limite retorna o pedido a neutro; para inverter, solicitar neutro primeiro. Solicitar parada trava o pedido e zera a marcha solicitada; liberar não retoma a marcha anterior. O estado e sua revisão vivem na RAM do mestre e reiniciam com ele.

Esses campos \textbf{não alteram a marcha publicada, os sensores, GPIO, PWM, DAC, relé, freio ou motor}. O campo de atuação habilitada permanece falso. A confirmação de parada registra um pedido, não comprova desaceleração ou imobilização. O botão da página é enviado sem diálogo de confirmação, mas ainda depende de navegador, Wi-Fi, UART e mestre; ele não substitui a parada de emergência física. Dois navegadores podem modificar o mesmo estado; a revisão permite perceber mudança, mas não fornece exclusividade de operador. Antes de qualquer integração física, são necessários intertravamentos baseados na velocidade medida, autorização individual, validade temporal dos pedidos, política para perda de enlace e retorno de estado do atuador.

\section{Terminal serial e procedimentos rápidos}
No terminal do Mestre a 115\,200 baud, a opção 1 solicita resumo de diagnóstico; 2 abre OTA; 3 e 4 tratam intervalos e pausa de relatórios; 5 gerencia redes; 6 consulta satélites; 7 mostra sensores, odômetro e falhas; 8 inicia envio dos CSVs pendentes; 9 grava áudio WAV temporário para treinamento; 10 ajusta avisos, logs e cores do menu; 11 mostra um estado rápido; 0 redesenha o menu. No submenu de sensores é possível pausar cada fonte ou todas para depuração e retomá-las sem recompilar. A pausa não deve ser confundida com falha do sensor; conferir o estado antes de usar a telemetria. A opção de logs permite escolher o nível global ou por serviço, incluindo silenciar, só erros, avisos e informações. Essas escolhas controlam o terminal; falhas relevantes continuam destinadas ao registro diagnóstico no SD quando disponível. A senha de uma rede aberta é vazia. A digitação de senhas Wi-Fi não ecoa caracteres; a senha de permissão de mudanças protegidas é mascarada. Para uma alteração, registrar estado antes e depois, resposta, versão dos três firmwares e log de erro caso o resultado não corresponda ao pedido.

\chapter{Wi-Fi, painéis web e acesso à telemetria}\label{cap:wifi-paineis}

\section{Função de cada rádio}
Os três ESP32-S3 não formam uma única rede IP. O ESP Equipe cria um ponto de acesso (\emph{Access Point}, AP) para operação e diagnóstico local. O ESP Visitantes cria outro AP, dedicado à apresentação pública da telemetria. O ESP Mestre atua como estação (\emph{STA}) quando precisa alcançar uma rede externa para atualização pela rede (\emph{Over-the-Air}, OTA) ou sincronização com o Google. No uso local, o Mestre troca dados com os satélites pela UART; os navegadores não acessam diretamente o Mestre por IP.

Essa separação mantém o painel público sem caminho HTTP para comandos, preserva a interface da Equipe quando o roteador externo não está disponível e dispensa Internet para visualizar as amostras. A rede externa do Mestre e os dois APs podem coexistir, mas alcance e desempenho dependem da alimentação, antenas, canais e ambiente de rádio. A arquitetura de UART e os quadros são descritos em outros capítulos; aqui o foco é como os dados chegam ao navegador e como se controla o acesso.

\begin{figure}[htbp]\centering
\begin{tikzpicture}[node distance=8mm and 11mm, >={Latex},
    bloco/.style={draw=nexus,rounded corners,fill=pale,align=center,font=\small,minimum height=10mm,text width=2.35cm}]
\node[bloco] (mestre) {ESP Mestre\\STA sob demanda};
\node[bloco,above right=9mm and 16mm of mestre] (equipe) {ESP Equipe\\AP e HTTP};
\node[bloco,below right=9mm and 16mm of mestre] (visitantes) {ESP Visitantes\\AP e HTTP};
\node[bloco,right=16mm of equipe] (operador) {Navegador\\da Equipe};
\node[bloco,right=16mm of visitantes] (publico) {Navegador\\do visitante};
\node[bloco,below=22mm of mestre] (roteador) {Roteador externo\\OTA e Google};
\draw[<->,thick] (mestre.north east)--node[above,sloped,font=\scriptsize]{UART1} (equipe.south west);
\draw[<->,thick] (mestre.south east)--node[below,sloped,font=\scriptsize]{UART2} (visitantes.north west);
\draw[<->,thick] (equipe)--node[above,font=\scriptsize]{Wi-Fi/HTTP} (operador);
\draw[<->,thick] (visitantes)--node[below,font=\scriptsize]{Wi-Fi/HTTP} (publico);
\draw[<->,thick,dashed] (roteador)--node[left,font=\scriptsize]{Wi-Fi STA} (mestre);
\end{tikzpicture}
\caption{Caminhos implementados. O Mestre não é um roteador entre os dois APs.}\label{fig:fluxo-web}
\end{figure}

\section{Configuração dos APs locais}\label{sec:aps-locais}
Os nomes e parâmetros a seguir constam das configurações dos respectivos firmwares. Na revisão consultada, os dois APs usam a mesma senha configurada no firmware; ela não é reproduzida neste manual. Ambos usam WPA2-PSK e servidor HTTP na porta 80. \emph{PSK} é a chave compartilhada usada para ingressar na rede. A senha comum simplifica a distribuição inicial, mas amplia o alcance de uma credencial divulgada: publicar o QR de Visitantes também permite a associação ao AP Equipe. \textbf{Não há chave administrativa HTTP adicional no painel da Equipe}: quem consegue acessar esse AP pode enviar comandos. Restringir fisicamente e por distribuição de credenciais o acesso à rede técnica. Se a política exigir PSKs diferentes, será preciso alterar a configuração, compilar ambos os firmwares e rever os QR.

\begin{table}[htbp]\centering\small
\caption{Configuração local dos dois pontos de acesso}\label{tab:aps-wifi}
\begin{tabularx}{\textwidth}{@{}>{\raggedright\arraybackslash\bfseries}p{3.2cm}X X@{}}
\toprule
Parâmetro & Equipe & Visitantes\\
\midrule
Nome da rede (SSID) & \texttt{NEXUS-Equipe} & \texttt{NEXUS-Visitantes}\\
Canal em 2,4 GHz & 11 & 1\\
Limite de clientes conectados & 5 & 10\\
Endereço para abrir o painel & \url{http://192.168.4.1/} & \url{http://192.168.4.1/}\\
Tempo padrão por conexão & 15 minutos; configurável & 15 minutos; configurável\\
Finalidade da interface & Diagnóstico e comandos para clientes do AP & Telemetria pública, somente leitura\\
\bottomrule
\end{tabularx}
\end{table}

Ambos usam WPA2-PSK e fornecem endereços por DHCP na própria rede. A UART1 de cada satélite liga-se exclusivamente ao mestre; sua pinagem e velocidade pertencem ao enlace interno, não à configuração do AP. No Visitantes, o firmware solicita largura de 20 MHz e potência máxima de 16 dBm. No Equipe, esses dois parâmetros seguem a configuração do driver; a potência aplicada deve ser lida no diagnóstico do dispositivo.

\emph{SSID} é o nome anunciado pela rede. Os canais 1 e 11 afastam os APs na faixa de 2,4 GHz e podem reduzir sobreposição; essa justificativa é uma inferência de engenharia, não uma medição de interferência no veículo. O limite maior de clientes do Visitantes atende ao público; o Equipe reserva sua interface de operação para um grupo menor. Os limites são de associações simultâneas, não de requisições HTTP por segundo.

O mesmo IP \texttt{192.168.4.1} nos dois APs não cria conflito enquanto forem redes distintas. O telefone ou computador recebe um endereço por \emph{Dynamic Host Configuration Protocol} (DHCP) ao associar-se a um AP e envia pacotes ao ESP daquela rede. Um dispositivo conectado aos dois SSIDs por interfaces diferentes deve ter suas rotas verificadas; normalmente, conecte-se apenas ao SSID pretendido. O sistema operacional pode avisar que a rede local está sem Internet: mantenha a conexão para abrir o painel.

\subsection{Portal cativo}
O Equipe anuncia o endereço do portal pela opção 114 do DHCP quando o cliente a suporta. Ele não mantém um DNS local para interceptar qualquer nome digitado. O Visitantes oferece um DNS local opcional na porta 53, que responde com o IP do AP para favorecer a detecção de portal cativo; se a tarefa DNS falhar, o IP direto ainda funciona. Requisições HTTP desconhecidas no Visitantes são redirecionadas ao painel. A abertura automática depende do sistema operacional e da política do telefone. Em ambos os casos a entrada confiável é \url{http://192.168.4.1/}.

\subsection{Tempo de sessão dos clientes}
Cada AP acompanha o tempo de associação. O limite padrão é 15 minutos, salvo em \emph{Non-Volatile Storage} (NVS) no espaço \texttt{wifi\_sessao}, chave \texttt{limite\_min}. NVS é a área persistente da flash separada do programa. O valor 0 desativa o limite; o máximo aceito é 1440 minutos. Ao vencer o prazo, o AP desconecta a estação. Isso libera uma vaga, mas não revoga a senha, não altera a concessão DHCP e não impede reconexão com a PSK salva. Para retirar o acesso, é preciso alterar a senha Wi-Fi nos firmwares e substituir os QR divulgados.

\section{Rede externa do Mestre}\label{sec:wifi-mestre}
O Mestre usa Wi-Fi em modo STA para alcançar um roteador externo. O rádio fica normalmente desligado e é acionado para operações que exigem rede, como OTA ou sincronização com o Google. Isso reduz consumo e atividade de rádio durante a instrumentação. Estar conectado ao AP Equipe não fornece Internet ao Mestre: o caminho de controle é UART e os satélites não são roteadores.

O firmware prevê uma rede principal de compilação, até cinco secundárias salvas em NVS e uma rede temporária preparada para teste. SSID e senha da principal ficam em cabeçalho privado \path{firmware/mestre-s3/include/nucleo/segredos_ota.h}. As secundárias e a política persistem no namespace \texttt{wifi\_perfis}. A principal também pode ser registrada em NVS para sobreviver a uma OTA cujo binário não contenha o mesmo cabeçalho privado. Apagar a NVS ou trocar a placa exige inventariar previamente os perfis.

\begin{table}[htbp]\centering\small
\begin{tabularx}{\textwidth}{>{\bfseries}p{5.5cm}X}
\toprule
Parâmetro padrão & Valor e função\\
\midrule
Canal preferencial & 0: procurar qualquer canal; o canal efetivo é imposto pelo roteador na associação\\
Largura e potência solicitadas & 20 MHz; máximo de 16 dBm, sujeito ao rádio e à região\\
Tentativas por rede & 3; faixa configurável de 1 a 10\\
Tempo por tentativa & 8000 ms; faixa de 2000 a 30000 ms\\
Intervalo entre tentativas & 750 ms; faixa de 100 a 5000 ms\\
Prazo máximo da sessão & 30000 ms; faixa de 5000 a 180000 ms, sempre ao menos igual ao tempo de uma tentativa\\
Redes secundárias & até 5 perfis persistentes\\
\bottomrule
\end{tabularx}
\caption{Política padrão em \path{configuracao_wifi.h}. O prazo de sessão do Mestre é distinto dos 15 minutos por cliente dos APs.}\label{tab:wifi-mestre}
\end{table}

O prazo máximo impede que uma rede ausente paralise indefinidamente a operação. Antes de associar-se, o Mestre varre as redes visíveis e tenta primeiro os perfis conhecidos que aparecem no resultado, respeitando a política de autorização das secundárias. Se a varredura falhar, volta à ordem normal dos perfis; se funcionar e não encontrar nenhuma rede elegível, registra esse motivo sem gastar todas as tentativas de associação. Uma rede oculta pode requerer investigação específica porque não aparece nessa seleção. Três tentativas curtas, separadas por 750 ms, acomodam falhas transitórias. A sessão pode acabar antes de esgotar todas as tentativas. O firmware serializa o rádio entre serviços; uma ação recebida no Equipe atravessa a UART e pode demorar enquanto o Mestre tenta associar-se.

\subsection{Perfis, autorização e procedimento}\label{sec:perfis-wifi}
\begin{enumerate}
  \item Definir a principal no cabeçalho privado antes da compilação, ou conservar a principal persistida. A geração inicial pode usar as variáveis \texttt{OTA\_WIFI\_SSID} e \texttt{OTA\_WIFI\_SENHA} com \path{ferramentas/gerar_segredos_ota.py}. O gerador reescreve o cabeçalho e inclui apenas as duas macros Wi-Fi; se já houver URL/token Google ou senha de permissão, editar manualmente para preservá-los.
  \item No painel Equipe, preparar temporariamente SSID e senha de uma rede externa e testar a conexão. Deixar a senha vazia para uma rede aberta; nesse caso, confirmar também que o roteador anuncia a rede sem proteção. A rede preparada ainda não ocupa uma das cinco posições persistentes.
  \item Salvar como secundária ou, com a senha de permissão do Mestre, promover a preparada a principal. Também é possível selecionar uma rede salva por índice de 1 a 5, excluir um perfil e restaurar a principal de compilação.
  \item Ajustar, se necessário, busca automática e limites de tentativas, duração e sessão. A ativação automática de secundárias e mudanças protegidas exigem a senha de permissão do Mestre.
  \item Depois de \emph{Conectar}, acompanhar o estado do Wi-Fi até ver associação ou erro. Depois de \emph{Salvar}, consultar a lista de redes salvas e conferir SSID e posição do perfil; conexão bem-sucedida não significa gravação persistente. O aceite do comando tampouco confirma a conclusão.
\end{enumerate}
Em uma rede usada só na manutenção, a autorização manual reduz associações inesperadas. O modo automático pode procurar alternativas conhecidas, mas a ordem dos perfis deve ser documentada. Alterar a principal sem testar a preparada pode deixar OTA e Google sem rota. Conectar ao roteador não prova que DNS, Internet, servidor OTA ou credenciais Google estejam funcionais.

\section{Acesso local e permissão do Mestre}\label{sec:acesso-web}
O AP Equipe exige a senha WPA2 para associação. A API HTTP local não pede uma segunda chave administrativa para consultas ou ações. Assim, um cliente admitido na rede técnica pode enviar pedidos ao satélite, inclusive comandos que este encaminha ao Mestre. O painel Visitantes continua separado e sem rotas de comando. Não publicar o QR ou a senha do AP Equipe como material público; o HTTP local não usa TLS.

Determinadas mudanças de perfis e política Wi-Fi conservam uma senha de permissão verificada pelo \emph{Mestre}. Essa proteção é específica dessas operações e não protege os demais comandos do painel. Para senhas de redes externas, o campo pode ficar vazio somente quando a rede for aberta. No terminal do Mestre, a entrada de senha Wi-Fi não ecoa caracteres; a senha de permissão de comandos protegidos usa máscara. Não registrar credenciais em capturas ou logs compartilhados.

\begin{figure}[htbp]\centering
\begin{tikzpicture}[node distance=11mm, >={Latex},
  etapa/.style={draw=nexus,rounded corners,fill=pale,align=center,text width=3.15cm,font=\small,minimum height=13mm}]
\node[etapa] (wifi) {1. Associação ao AP\\PSK WPA2};
\node[etapa,right=of wifi] (uart) {2. Pedido HTTP\\encaminhado pela UART};
\node[etapa,right=of uart] (mestre) {3. Se exigida,\\permissão do Mestre};
\draw[->,thick] (wifi)--(uart);
\draw[->,thick] (uart)--(mestre);
\end{tikzpicture}
\caption{Caminho de uma ação técnica. A permissão adicional só se aplica às mudanças protegidas de Wi-Fi; a resposta HTTP não confirma o término da ação.}\label{fig:camadas-acesso}
\end{figure}

\section{Painel da Equipe}\label{sec:painel-equipe}
O servidor entrega a mesma aplicação web em \texttt{/}, \texttt{/painel}, \texttt{/diagnostico} e \texttt{/mestre}. As sete páginas têm funções diferentes:
\begin{description}
  \item[Telemetria:] velocidade, movimento, energia quando disponível, sessão atual, qualidade do enlace e clientes dos APs. É a visão rápida do estado e da idade do dado.
  \item[Comandos:] operações de Google, simulador, rede externa do Mestre, clientes do AP e OTA. Algumas mudanças de Wi-Fi exigem a permissão do Mestre.
  \item[Controle:] histórico recente e campos de marcha solicitada, limites e parada. A página apresenta dados e solicitações; não fecha um laço de atuação física do veículo.
  \item[ESP Mestre:] versão, saúde, memória, temperatura, tarefas, sensores, Wi-Fi, UART, sincronização, OTA e configuração conforme as respostas recebidas pela UART.
  \item[ESP Equipe:] diagnóstico local: versão, Wi-Fi, memória, CPU, tarefas, temperatura, UART, comandos e OTA.
  \item[ESP Visitantes:] estado do satélite público, versão e enlace alcançável pelo Mestre/Equipe. Não é um caminho de comandos para o público.
  \item[Sobre:] créditos, identidade visual e contexto do sistema.
\end{description}

Na página \emph{Telemetria} da Equipe, as leituras do INA228 ficam dentro do cartão \emph{Operação e sensores}, junto das demais informações de aquisição, sem um cartão INA separado. O cartão \emph{Bateria e consumo} mostra pelo menos duas casas decimais nos valores disponíveis e três na tensão; a formatação não aumenta a resolução da medição. O diagnóstico do ESP Mestre também permite verificar presença, validade e erros do sensor. Tensão, corrente, potência, queda no shunt, temperatura e acumuladores só devem ser interpretados com a origem e o estado de leitura correspondentes; ausência do módulo ou sensor pausado não equivale a valor zero.

Os arquivos de origem são \path{firmware/equipe-s3/web/index.html}, \path{style.css}, \path{layout.css} e \path{app.js}. O script \path{web/gerar_painel.py} os embute em \path{src/servicos/painel_web.c} antes da compilação. Editar as fontes e regenerar o C; uma mudança direta no arquivo gerado se perderá. O logotipo comum vem de \path{firmware/compartilhado/web/logo_carro_eletrico.png}.

O navegador busca telemetria a 1, 2, 3 ou 4 Hz, sendo 4 Hz o padrão (uma tentativa a cada 250 ms). Estados de OTA, Google e simulador são atualizados em torno de 2 s, independentemente da cadência do gráfico. Latência HTTP, contenda do rádio e resposta lenta podem reduzir a taxa efetiva. Histórico e exportação CSV são montados no navegador somente enquanto a página permanece aberta; recarregar ou fechar a aba encerra esse histórico local. Isso não equivale a armazenamento de bordo de longo prazo.

Ler cada campo junto de validade, idade e origem. Ausência de amostra não deve ser convertida em zero. O Mestre já lê o INA228 por I\textsuperscript{2}C para obter tensão do barramento, queda no shunt, corrente, potência, temperatura interna do circuito e acumuladores de energia/carga, quando o módulo responde e a medição é válida. O shunt configurado é de 1\,m$\Omega$ e a corrente máxima prevista é 75\,A; conferir montagem, polaridade e calibração antes de tratar valores como ensaio validado. Os sinais físicos de acelerador e freio ainda não são adquiridos. O simulador gera dados marcados como simulados; uma grandeza que só existe nele não deve ser apresentada como medição real.

No gráfico do sistema elétrico, as séries principais usam as escalas laterais esquerda e direita. Uma terceira série pode ser ligada opcionalmente e começa desligada para preservar a leitura. Sua escala própria aparece na legenda, e passar o mouse ou tocar o traçado revela o valor da amostra. Potência, corrente, tensão e estado de carga (SoC) são escolhas quando há dados válidos; autonomia restante é uma estimativa disponível apenas na simulação. Queda no shunt, temperatura do INA228, energia e carga acumuladas aparecem somente para medições reais válidas. Misturar grandezas no mesmo gráfico permite comparar tendências, mas cada unidade e escala devem ser lidas separadamente.

\subsection{API HTTP do Equipe}\label{sec:api-equipe}
\begin{table}[htbp]\centering\small
\begin{tabularx}{\textwidth}{>{\ttfamily}p{4.5cm}p{1.4cm}X}
\toprule
\normalfont\bfseries Rota & \bfseries Método & \bfseries Finalidade\\
\midrule
/api/telemetria & GET & Última amostra, sequência, idade e métricas UART.\\
/api/diagnostico & GET & Versão, rádio, memória, CPU/tarefas, temperatura, UART, comandos e OTA locais.\\
/api/wifi/clientes & GET & Estações associadas, MAC, tempo de conexão e limite de sessão.\\
/api/wifi/limite & POST & Alterar limite dos clientes do AP.\\
/api/wifi/desconectar & POST & Desconectar cliente do AP.\\
/api/mestre?consulta=... & GET & Consultar o Mestre pela UART e devolver JSON.\\
/api/mestre?consulta=... & POST & Encaminhar ação ao Mestre; a conclusão é assíncrona.\\
\bottomrule
\end{tabularx}
\caption{Rotas principais implementadas. Não há autenticação HTTP adicional; certas ações de Wi-Fi têm permissão adicional no Mestre.}\label{tab:apis-web}
\end{table}

As consultas GET do encaminhamento incluem \texttt{ping}, \texttt{versao}, \texttt{saude}, \texttt{telemetria}, \texttt{uart}, \texttt{ota}, \texttt{sincronizacao}, \texttt{wifi}, \texttt{diagnostico\_mestre}, \texttt{configuracao}, \texttt{simulador}, \texttt{rede\_wifi\_salva} com índice de 1 a 5, além de \texttt{ping} e \texttt{versao} para cada satélite. Uma consulta inválida ou método inadequado recebe HTTP 400/405, conforme o caso.

As ações POST abrangem \texttt{verificar\_ota}, \texttt{autorizar\_mestre}, \texttt{autorizar\_equipe}, \texttt{autorizar\_visitantes}, \texttt{autorizar\_todos}, \texttt{cancelar\_ota}, \texttt{resumo\_terminal}, \texttt{alternar\_relatorios}, \texttt{zerar\_percurso}, \texttt{sincronizar\_google}, \texttt{intervalo\_relatorio}, \texttt{intervalo\_ota}, \texttt{intervalo\_telemetria\_ota}, \texttt{simulador\_telemetria}, \texttt{wifi\_preparar}, \texttt{wifi\_conectar\_preparada}, \texttt{wifi\_salvar\_preparada}, \texttt{wifi\_conectar\_salva}, \texttt{wifi\_excluir\_salva}, \texttt{wifi\_automatico}, \texttt{wifi\_alterar\_principal}, \texttt{wifi\_restaurar\_principal} e \texttt{wifi\_politica}. A carga exata e os limites devem ser conferidos no despachante e em \path{documentacao/satelites.md} antes de automatizar chamadas.

Para \texttt{wifi\_preparar}, o corpo informa SSID e senha em linhas separadas. Ações sobre perfil salvo usam índice de 1 a 5. Alternâncias, como simulador e busca automática, usam 0/1; ações protegidas acrescentam a permissão do Mestre. \texttt{wifi\_politica} transmite quantidade de tentativas, duração por tentativa, intervalo, sessão máxima e permissão, nessa ordem. O formulário do painel reduz erros de formato. A resposta \texttt{ok:true} de um POST significa que o pedido foi aceito, não que rede, sincronização ou OTA terminou. Consultar \texttt{wifi}, \texttt{sincronizacao} ou \texttt{ota} depois da ação. Uma conexão Wi-Fi pode levar aproximadamente 185 s em cenário de esgotamento de tentativas.

O servidor admite no máximo duas consultas simultâneas ao Mestre e uma ação. Excesso pode receber HTTP 503 (\emph{Service Unavailable}); aguardar e repetir uma consulta, sem reenviar cegamente uma ação talvez já aceita. Erros de UART, idade da amostra e falha de conexão externa precisam ser lidos nos diagnósticos próprios, não deduzidos apenas de um cartão vazio.

\section{Painel Visitantes e validade da amostra}\label{sec:painel-visitantes}
O Visitantes oferece somente \texttt{GET /} e \texttt{GET /api/telemetria}. Não há página pública de comandos, diagnóstico detalhado ou gerenciamento da rede externa. O público recebe velocidade, movimento, gráficos e energia estimada apenas quando os campos necessários existem. O velocímetro usa escala visual fixa de 0 a 35 km/h; essa escala não define limite do veículo. O histórico gráfico começa quando a aba é aberta.

O navegador consulta a API a cada 500 ms (2 Hz); o Mestre publica telemetria aos satélites a 4 Hz. Nem toda publicação vira um ponto no gráfico. O JSON inclui \texttt{valid}, idade, sequência e marca de simulação. Um campo indisponível é \texttt{null}, não zero. Após mais de 2,5 s sem amostra nova, a página insere lacuna no gráfico para não inventar continuidade durante perda de UART ou parada do emissor. Distinguir último valor recebido, valor válido e valor atual ao interpretar uma captura.

\section{Manutenção e diagnóstico dos painéis}\label{sec:qr-paineis}
Os PNG/SVG em \path{artefatos/qr_visitantes} incluem um QR de associação Wi-Fi e outro com a URL do painel. O primeiro contém a senha do AP e deve ser tratado como material de acesso, inclusive em fotos e impressões. Se SSID, senha ou IP mudar, executar \path{ferramentas/gerar_qr_visitantes.py}, validar com um telefone e substituir impressos antigos. Regenerar um QR não altera o ESP: também é preciso compilar, gravar e verificar o firmware.

Para revisar a interface no computador, há \path{ferramentas/visualizar_painel_equipe.py} (porta 8766 com \texttt{--rede}) e \path{ferramentas/visualizar_painel_visitantes.py} (porta 8765). Eles usam APIs locais simuladas. A prévia testa layout e navegação, mas não valida UART, temporização real, alcance Wi-Fi nem sensores. Depois de editar as fontes web, executar o gerador, compilar, gravar e testar pelo AP físico. O serial confirma SSID/canal, IP, falhas do DNS, potência efetiva e início do HTTP.

\subsection{Diagnóstico por sintoma}\label{sec:diagnostico-wifi}
\begin{table}[htbp]\centering\small
\begin{tabularx}{\textwidth}{>{\bfseries}p{4.4cm}X}
\toprule
Sintoma & Sequência de verificação\\
\midrule
SSID não aparece & Confirmar alimentação e serial; conferir firmware do satélite correto, antena, distância, canal e falha do driver Wi-Fi.\\
Conecta, mas o portal não abre & Permanecer no AP apesar do aviso de ausência de Internet; abrir \url{http://192.168.4.1/}; conferir DHCP e rota da interface ativa.\\
Rede externa sem senha falha & Confirmar que a rede é realmente aberta e deixar a senha vazia; consultar o estado Wi-Fi do Mestre e a mensagem de erro.\\
Equipe não abre após OTA & Verificar serial, inicialização do AP/HTTP, autoexame e possível rollback; depois IP do AP.\\
HTTP 503 ou ação sem resultado & Reduzir consultas concorrentes, consultar estado após POST e verificar UART. Evitar repetir ação apenas porque a página demorou.\\
Gráfico parado ou com lacuna & Conferir idade, sequência e validade; verificar UART e publicação do Mestre. No Visitantes, 2,5 s sem amostra gera lacuna intencional.\\
OTA/Google não conectam & Consultar perfil e política do Mestre, autorização, associação, prazo, DNS/Internet e credenciais do serviço. O AP local não fornece essa rota.\\
Cliente cai após 15 min & Conferir \texttt{limite\_min}; é a política padrão do AP e o cliente pode se associar de novo.\\
\bottomrule
\end{tabularx}
\caption{Primeiros passos. A interpretação final exige registros dos ESPs e confirmação no hardware montado.}\label{tab:diagnostico-wifi}
\end{table}

\paragraph{Limites desta configuração.}
Este capítulo descreve os valores e rotas identificados no firmware e na documentação do repositório. A escolha dos canais como forma de reduzir interferência é uma inferência técnica a validar no local. A abertura automática do portal, potência efetiva, comportamento com muitos telefones, alcance e conclusão de OTA/Google em uma rede específica exigem ensaio no hardware montado. A integração de código do INA228 não substitui conferência da instalação elétrica e calibração com instrumento de referência.

\chapter{Memórias, histórico e integração com o Google}\label{cap:memorias_google}

\section{Por que há várias memórias}
O NEXUS não possui uma memória única que sirva igualmente para executar código, atender periféricos, guardar ajustes e conservar horas de telemetria. O mestre usa um ESP32-S3-WROOM-1-N16R8, isto é, módulo previsto para 16\,MiB de flash e 8\,MiB de PSRAM. A flash é persistente, porém tem ciclos limitados de escrita; a RAM interna e a PSRAM são voláteis; o cartão microSD é removível e comporta arquivos longos. Os satélites usam uma tabela de partições para flash de 4\,MiB. Antes de gravar a placa, é obrigatório conferir o módulo instalado: a designação do ambiente de compilação não altera a capacidade física.

\begin{table}[htbp]\centering\small
\caption{Papel e persistência das memórias no sistema}\label{tab:memorias_papel}
\begin{tabularx}{\textwidth}{@{}p{2.5cm}X X@{}}\toprule
Meio & Uso no projeto e motivo & Após perda de energia\\\midrule
Flash de programa & Guarda as duas imagens de aplicação OTA e os dados de partida; permite voltar a uma imagem anterior. & Permanece, salvo apagamento ou regravação.\\
NVS na flash & Pequenos valores alteráveis: odômetro, perfis e política Wi-Fi, intervalo OTA e ajustes específicos dos satélites. Evita regravar uma imagem inteira para mudar uma preferência. & Permanece enquanto a partição NVS não for apagada nem corrompida.\\
RAM interna & Tarefas, pilhas, filas e buffers que exigem acesso rápido. Parte com capacidade DMA é necessária para transferências de periféricos. & Perdida.\\
PSRAM externa & Históricos e buffers grandes que não precisam ocupar a RAM interna escassa. O firmware mede disponibilidade e admite alguns recuos para RAM interna. & Perdida.\\
RTC & Domínio de memória do chip que pode sobreviver a certos tipos de reinício, conforme a implementação; não substitui NVS ou cartão. & Não se deve presumir retenção após corte de alimentação.\\
microSD & CSVs, logs e cópia de configuração. Permite retirar e auditar arquivos sem regravar o ESP. & Permanece se o cartão e o sistema de arquivos estiverem íntegros.\\\bottomrule
\end{tabularx}
\end{table}

\begin{nota}{Estado implementado versus propriedade do chip}
A arquitetura usa concretamente flash, NVS, RAM interna, PSRAM e microSD. A memória RTC aparece aqui para evitar confundi-la com armazenamento durável: nos serviços de histórico e configuração analisados não há dependência de uma cópia em RTC. Em particular, o estado de seleção das imagens OTA pertence aos metadados de OTA na flash, não a um arquivo CSV nem a uma suposta cópia permanente em RAM.
\end{nota}

\subsection{Mapa da flash e consequências para OTA}\label{sec:particoes_flash}
A tabela de partições fixa regiões e limites de tamanho antes da compilação. \texttt{ota\_0} e \texttt{ota\_1} são espaços alternativos para executáveis; uma atualização escreve a partição inativa e a seleção de partida muda somente segundo o fluxo OTA. \texttt{otadata} guarda os metadados de escolha e validação da imagem; \texttt{nvs} guarda valores de aplicação; \texttt{phy\_init} é reservada ao subsistema de rádio. Nenhuma dessas partições é uma área geral para CSV. Um firmware maior que sua partição OTA não pode ser instalado integralmente nesse esquema.

\begin{table}[htbp]\centering\small
\caption{Partições configuradas nos projetos PlatformIO; valores em hexadecimal}\label{tab:particoes_nexus}
\begin{tabularx}{\textwidth}{@{}lXXl@{}}\toprule
Partição & Mestre, 16\,MiB & Cada satélite, 4\,MiB & Finalidade\\\midrule
\texttt{nvs} & \texttt{0x9000}, \texttt{0x6000} (24\,KiB) & \texttt{0x9000}, \texttt{0x5000} (20\,KiB) & Preferências persistentes\\
\texttt{otadata} & \texttt{0xF000}, \texttt{0x2000} (8\,KiB) & \texttt{0xE000}, \texttt{0x2000} (8\,KiB) & Seleção e estado OTA\\
\texttt{phy\_init} & \texttt{0x11000}, \texttt{0x1000} (4\,KiB) & \texttt{0x10000}, \texttt{0x1000} (4\,KiB) & Dados de rádio\\
\texttt{ota\_0} & \texttt{0x20000}, \texttt{0x600000} (6\,MiB) & \texttt{0x20000}, \texttt{0x1E0000} (1920\,KiB) & Imagem A\\
\texttt{ota\_1} & \texttt{0x620000}, \texttt{0x600000} (6\,MiB) & \texttt{0x200000}, \texttt{0x1E0000} (1920\,KiB) & Imagem B\\\bottomrule
\end{tabularx}
\end{table}

Os deslocamentos acima vêm de \path{firmware/mestre-s3/particoes_ota.csv} e das tabelas equivalentes de Equipe e Visitantes. Regiões não listadas incluem áreas convencionais de partida e espaços reservados; espaço livre no chip não deve ser tratado como partição utilizável sem alterar a tabela, reconstruir e validar a gravação. A atualização de firmware não deve apagar a NVS, pois isso perderia odômetro e redes; uma gravação que apague toda a flash tem consequência diferente. O mecanismo de confirmação ou retorno de uma nova imagem é explicado no capítulo de OTA.

\subsection{RAM, PSRAM, pilhas e transferências}
A RAM interna é o recurso mais restrito durante Wi-Fi e HTTPS: as tarefas têm pilhas próprias, os controladores precisam de buffers e uma alocação grande depende não só da soma de bytes livres, mas do \emph{maior bloco livre}. Por isso o diagnóstico publica memória interna livre, menor nível já observado, maior bloco contíguo e menor margem de pilha. Um total livre aparentemente suficiente ainda pode falhar ao reservar um bloco contíguo. Buffers destinados a DMA (\emph{Direct Memory Access}, transferência direta entre periférico e memória) precisam de capacidades de alocação compatíveis; PSRAM livre não equivale automaticamente a RAM interna adequada a DMA.

O mestre espera 8\,MiB de PSRAM. O histórico de saúde mantém 120 amostras a 1\,Hz nela; quando não está disponível, tenta uma versão reduzida de 30 amostras na RAM interna. Outros caminhos, como partes do pipeline OTA, também tentam PSRAM e registram quando recuam para RAM interna. Essa política permite continuar funções essenciais, mas diminui a margem de memória; não comprova que toda carga simultânea foi validada sem PSRAM. O registro CSV e o upload executam tarefas separadas no núcleo 0, enquanto aquisição e enlaces prioritários ficam configurados no núcleo 1. Essa distribuição reduz disputa, mas não transforma escrita em SD, rede e alocação em operações de tempo determinístico.

\section{Caminho de uma amostra até o arquivo}\label{sec:csv_fluxo}
O retrato ao vivo é produzido a 4\,Hz. A tarefa de histórico seleciona uma amostra por segundo, coloca-a em uma fila de 24 posições e deixa uma única tarefa proprietária da escrita de CSV. A fila absorve atrasos curtos do cartão; se a escrita ficar atrasada por mais tempo, amostras podem ser descartadas e o contador de perdas deve ser examinado. A separação impede que o tempo de resposta do cartão bloqueie diretamente o ciclo de aquisição. Os gráficos abertos no navegador contêm apenas o que aquela página recebeu e não são substitutos do arquivo local.

\begin{figure}[htbp]\centering
\begin{tikzpicture}[
  >=Latex, node distance=4mm,
  etapa/.style={draw=nexus,rounded corners,fill=pale,align=center,text width=10.8cm,minimum height=7mm,inner sep=3pt,font=\small},
  seta/.style={->,thick,draw=accent}
]
\node[etapa] (origem) {Retrato do mestre a 4\,Hz: aquisição, cálculo e indicação de validade};
\node[etapa,below=of origem] (fila) {Amostragem a 1\,Hz $\rightarrow$ fila de 24 posições $\rightarrow$ único gravador};
\node[etapa,below=of fila] (ativo) {microSD: \texttt{.tmp} ativo; sincronização a cada 10\,s};
\node[etapa,below=of ativo] (pendente) {Fechamento por 1\,h, 8\,MiB ou pedido de envio $\rightarrow$ \texttt{.csv} pendente};
\node[etapa,below=of pendente] (google) {Envio manual HTTPS $\rightarrow$ confirmação exata do Web App};
\node[etapa,below=of google] (enviado) {Movimentação local para \texttt{enviados/}; cópia original no Drive e linhas na planilha};
\draw[seta] (origem)--(fila); \draw[seta] (fila)--(ativo);
\draw[seta] (ativo)--(pendente); \draw[seta] (pendente)--(google);
\draw[seta] (google)--(enviado);
\end{tikzpicture}
\caption{Fluxo do histórico; o aceite de um comando não é confirmação do destino.}\label{fig:fluxo_historico}
\end{figure}

\begin{table}[htbp]\centering\footnotesize
\caption{Estados observáveis de um CSV no cartão}\label{tab:csv_estados}
\begin{tabularx}{\textwidth}{@{}p{3.7cm}X@{}}\toprule
Nome ou local & Significado e ação segura\\\midrule
\path{telemetria_XXXXXXXX_NNNN.tmp} & Arquivo ativo. Não é enviado enquanto está aberto. \texttt{XXXXXXXX} identifica a sessão de partida e \texttt{NNNN}, a sequência do arquivo; nenhum dos dois é data civil.\\
\path{telemetria_XXXXXXXX_NNNN.csv} & Arquivo fechado, elegível para envio. Continua pendente após perda de rede ou confirmação insuficiente.\\
\path{telemetria/enviados/} & Cópia local preservada após confirmação do nome pelo Web App e movimentação bem-sucedida no cartão.\\
\path{*.incompleto} & Arquivo que falhou na recuperação e é retido para inspeção, sem envio automático.\\\bottomrule
\end{tabularx}
\end{table}

O arquivo gira ao atingir uma hora ou 8\,MiB, o que ocorrer primeiro. Isso limita o trabalho de recuperação de um arquivo interrompido e evita um CSV sem limite prático. O serviço sincroniza o arquivo ativo com o cartão aproximadamente a cada 10\,s; uma queda abrupta de energia ainda pode perder as linhas que não chegaram ao meio. Na partida, um \texttt{.tmp} com cabeçalho esperado e última linha completa é renomeado para \texttt{.csv}; se a estrutura não passar nessa verificação, ele é preservado como \texttt{.incompleto}. Não se deve renomear manualmente um arquivo danificado para forçar upload.

Cada linha contém \texttt{tempo\_sessao\_ms}, ligado ao tempo desde a partida do mestre, e \texttt{simulado}, além dos campos da telemetria. A existência de uma coluna ou o número zero não demonstra que haja um sensor físico correspondente. Para analisar ensaios, confronte indicadores de validade, estado de simulação, revisão do firmware e calibração. O nome de sessão, sozinho, não fornece horário civil.

\subsection{Logs estruturados e contexto de incidentes}
O mestre escreve logs próprios em \path{/sdcard/logs}, em segmentos numerados, separados da saída \texttt{ESP\_LOG} do terminal. Em regime normal, o resumo de saúde é registrado a cada 30\,s. Ao detectar falha, o serviço tenta descarregar o histórico recente de saúde com marcadores \texttt{HIST\_INICIO}, \texttt{HIST} e \texttt{HIST\_FIM}, uma vez por incidente; depois registra a cada 2\,s por 2\,min. Para falha grave, usa 1\,s por 5\,min. Uma nova falha grave pode abrir outra captura. Segmentos chegam a 8\,MiB, com limite configurado de 16\,MiB por sessão. Esses limites contêm o uso do cartão sem prometer retenção ilimitada.

Se o SD estiver indisponível, o primeiro lote de contexto pode permanecer pendente em RAM; essa memória é volátil e não constitui garantia após corte de energia. Incidentes posteriores sem capacidade de captura são contabilizados. O serviço tenta novamente o cartão após 30\,s. O relógio civil só passa a ser usado quando obtém uma referência UTC válida do cabeçalho \texttt{Date} de uma resposta HTTPS de OTA ou Google; antes disso, a linha deve ser interpretada pelo tempo relativo à partida. Não retroatribuir um horário UTC às amostras anteriores sem uma correlação de ensaio documentada. Logs estruturados não fazem parte do upload Google dos CSVs.

\subsection{Cópia de configuração, NVS e recuperação}\label{sec:backup_config}
A NVS é a fonte persistente normal para pequenas preferências. O odômetro acumulado é salvo em milímetros numa chave própria; perfis de Wi-Fi, política de seleção de rede e intervalo OTA possuem seus espaços de persistência. Quando a NVS não puder ser aberta, o odômetro pode continuar volátil, com erro diagnosticável. Cada gravação requer \texttt{commit}; atualizar a variável em RAM não equivale a gravá-la na flash.

O mestre também escreve \path{/sdcard/configuracao_atual.txt} aproximadamente a cada 60\,s. O conteúdo tem versão de formato, sequência e checksum para rejeitar truncamento ou alteração acidental. A escrita passa por arquivo temporário e substituição; a sequência só avança após gravação bem-sucedida. Se falhar, há nova tentativa após 5\,s. Na partida, o serviço compara a cópia principal e a temporária e escolhe a válida mais recente. O backup não é espelho completo da NVS: restaura odômetro \emph{somente se a chave correspondente estiver ausente}, e aplica seletivamente preferências Wi-Fi, intervalo OTA e ajustes de relatórios segundo a situação das chaves. Não deve sobrescrever um valor NVS válido apenas por existir um arquivo mais antigo.

Os nomes de redes podem aparecer na cópia, mas as senhas Wi-Fi não aparecem. O token e a URL do Google, definidos no cabeçalho privado de compilação, também não são recuperados desse arquivo. Depois de apagar NVS ou trocar a placa, recadastre credenciais e confira a política de rede; depois de reconstruir o firmware, confira a configuração privada Google. Uma queda antes do próximo backup pode deixar uma diferença de até aproximadamente um minuto de progresso nessa cópia, e a persistência NVS do odômetro tem sua própria cadência. Preserve cartão e NVS originais ao investigar divergência; não use o backup como prova de que todo valor em RAM foi salvo.

\begin{nota}{Recuperação não é sincronização automática de dois bancos}
O microSD protege alguns ajustes contra ausência de chaves NVS, mas não fornece restauração integral de senhas, segredos ou todos os estados transitórios. A política de recuperação é seletiva para evitar que uma cópia antiga reverta uma configuração válida.
\end{nota}

\section{Preparação da integração com Google Apps Script}\label{sec:google_config}
O envio foi deliberadamente colocado fora do ciclo em tempo real e começa por pedido do operador. O destino é uma pasta do Google Drive com os CSVs originais e uma planilha \texttt{telemetria\_carro} na mesma pasta. Cada sessão de partida recebe uma aba \texttt{sessao\_XXXXXXXX}; a coluna \texttt{arquivo\_origem} vincula cada linha ao CSV de origem. O ESP não executa OAuth interativo: o Web App roda com as permissões da conta que o publicou. Como o acesso do aplicativo precisa permitir requisições sem login Google no ESP, um token compartilhado longo constitui a autorização efetiva. URL e token devem permanecer fora do repositório público e de capturas de log.

\begin{enumerate}
\item Crie, na conta que conservará os dados, uma pasta no Drive e copie o \emph{ID} presente no endereço da pasta; o nome visível não é o ID.
\item Crie um projeto em Google Apps Script e coloque nele o conteúdo de \path{documentacao/telemetria-google/Code.gs}. Nas propriedades do script, crie \texttt{ID\_PASTA\_TELEMETRIA} com o ID da pasta e \texttt{TOKEN\_TELEMETRIA} com segredo aleatório de pelo menos 32 caracteres. A propriedade \texttt{ID\_PLANILHA\_TELEMETRIA} é preenchida pelo script ao criar ou localizar a planilha.
\item Implante como \emph{Aplicativo da Web}, executando como a própria conta e permitindo acesso a \emph{Qualquer pessoa}; uma configuração que exija login Google impede a requisição direta do ESP. Autorize Drive e Sheets e copie a URL publicada terminada em \texttt{/exec}. A URL \texttt{/dev} serve ao editor e não é a configuração de produção.
\item No mestre, crie o cabeçalho privado \path{firmware/mestre-s3/include/nucleo/segredos_ota.h} a partir do exemplo. Configure uma rede Wi-Fi externa com internet, \texttt{TELEMETRIA\_GOOGLE\_URL} com a URL \texttt{/exec} e \texttt{TELEMETRIA\_GOOGLE\_TOKEN} com o mesmo token das propriedades do script. Compile e grave o mestre.
\item Se modificar \texttt{Code.gs} depois, use \emph{Gerenciar implantações}, edite a implantação, selecione \emph{Nova versão} e implante de novo. Salvar o código no editor não atualiza por si só a versão \texttt{/exec} consumida pelo ESP.
\end{enumerate}

O script \path{ferramentas/gerar_segredos_ota.py} gera somente as macros iniciais de Wi-Fi e \emph{sobrescreve} o cabeçalho privado. Depois de adicionar URL e token Google, edite o cabeçalho diretamente para não perdê-los. O arquivo real é local e ignorado pelo Git; o manual não contém credenciais reais. O token deve ser tratado como segredo operacional: se for exposto, troque-o no Apps Script e no firmware mestre, reconstrua e grave uma nova versão.

\subsection{Protocolo de envio, confirmação e duplicatas}\label{sec:google_fluxo}
O operador pode solicitar sincronização pelo terminal serial do mestre a 115200\,baud ou pela página \emph{Comandos} do ESP da Equipe, depois de habilitar acesso administrativo. A página de Visitantes não tem esse comando. O mestre não inicia dois envios simultâneos nem toma uma sessão Wi-Fi já ocupada por outra operação. Ele fecha o CSV atual, abre uma sessão temporária da estação Wi-Fi na rede externa, enumera os \texttt{.csv} pendentes e faz um \texttt{POST} HTTPS por arquivo. O conteúdo é transferido em blocos de 8192 bytes; o pedido inclui os parâmetros \texttt{arquivo} e \texttt{token}. A configuração prevê 60\,s de limite por upload, até quatro tentativas por arquivo e esperas crescentes entre tentativas. O painel mostra arquivo atual, bytes, porcentagem, taxa e resultado, mas o aceite inicial do comando significa apenas que a tarefa começou.

O serviço segue o redirecionamento de resposta que o Google pode fazer para \texttt{script.googleusercontent.com}. A condição de sucesso é conjunta: código HTTP 200 \emph{e} corpo exatamente \texttt{OK:} seguido do mesmo nome de arquivo enviado. Somente depois disso o mestre move o arquivo para \path{/sdcard/telemetria/enviados}. Código 200 isolado, resposta parcial, página de login, \texttt{ERRO:*} ou confirmação de outro nome não são comprovantes. Se as quatro tentativas falharem, o CSV permanece pendente e o lote termina com erro. Se não existir arquivo pendente, a operação pode terminar com sucesso sem transferir bytes.

O Apps Script verifica antes de inserir linhas se aquele nome já consta da aba da sessão e verifica antes de criar um CSV se o mesmo nome já existe na pasta. Isso torna a repetição do mesmo arquivo segura nos casos previstos: por exemplo, o Google pode ter gravado o arquivo mas a resposta se perder na rede, ou o mestre pode receber \texttt{OK} e falhar ao mover o arquivo no SD. Na próxima tentativa o nome permanece igual, e o script evita duplicar arquivo e linhas. A deduplicação depende desse nome estável e do script instalado; não resolve alteração manual do nome, mudança de pasta ou planilhas antigas, que não são migradas automaticamente.

\begin{table}[htbp]\centering\small
\caption{Diagnóstico de envio e ação indicada}\label{tab:google_falhas}
\begin{tabularx}{\textwidth}{@{}p{4.0cm}X@{}}\toprule
Sinal & Verificação ou consequência\\\midrule
Pedido recusado ou botão desabilitado & Consultar estado do Mestre, microSD, Wi-Fi e sessão já ativa; não repetir comandos às cegas.\\
Sem pendentes & Confirmar contador de amostras e fechamento de CSV. Uma execução vazia não prova que houve envio anterior.\\
Erro de DNS, conexão ou TLS & Conferir perfil e sinal da rede externa, internet, URL \texttt{/exec}, data de resposta e registros do terminal. O AP local da Equipe não é, por si só, acesso à internet.\\
\texttt{ERRO:configuracao} ou \texttt{ERRO:autorizacao} & Conferir propriedades do script, ID de pasta e correspondência do token, sem imprimi-lo em relatórios.\\
\texttt{ERRO:nome}, \texttt{ERRO:conteudo} ou \texttt{ERRO:csv} & Inspecionar nome e conteúdo preservados no cartão; verificar cabeçalho e ao menos uma linha de dados.\\
\texttt{ERRO:interno}, HTTP inesperado & Examinar \emph{Execuções} do Apps Script, permissões da conta, pasta, implantação e resposta no terminal.\\
\texttt{OK} remoto, arquivo ainda pendente & Investigar falha de movimentação no SD. Reenviar o mesmo nome é seguro dentro da política de deduplicação.\\\bottomrule
\end{tabularx}
\end{table}

\subsection{Verificação em bancada e limites documentais}
Com cartão montado e uma rede externa configurada, deixe o mestre produzir amostras reais ou identificadas como simuladas. Confira que o diagnóstico mostra amostras gravadas, ausência de perdas na fila e arquivo fechado elegível. Solicite o envio, espere o \emph{resultado final}, confirme o CSV original na pasta, as linhas na aba da sessão e o arquivo em \texttt{enviados/}. Para testar recuperação, interrompa a rede antes da confirmação e verifique que o CSV continua pendente; ao restabelecer a rede, repita e confira que não apareceram duplicatas. Registre revisão do firmware, versão do script, nome do arquivo e observações, sem anotar o token.

O código implementa esse fluxo, mas a documentação consultada não constitui um relatório de ensaio exaustivo com todas as combinações de falha de energia, cartão removido, PSRAM ausente e rede intermitente. Na instalação física, devem ser verificados retenção após desligamento, taxas de perda em carga, capacidade do cartão, integridade da recuperação e comportamento de reenvio antes de usar o histórico como registro experimental definitivo.

\chapter{Atualização OTA, publicação e retorno de versão}\label{cap:ota}

\section{Objetivo e limites da atualização}
OTA significa \emph{Over-the-Air}: substituir o firmware sem repetir a gravação por cabo. No NEXUS, apenas o mestre obtém as imagens da Internet. Ele instala sua própria imagem por HTTPS ou transmite pela UART a imagem destinada a um satélite. A decisão de \emph{instalar} continua condicionada a uma autorização do operador. Uma versão encontrada, uma ordem aceita pelo servidor web, um bloco UART confirmado e uma imagem gravada são acontecimentos distintos; nenhum deles comprova, isoladamente, que o novo programa iniciou e permaneceu funcional.

Há três firmwares independentes: mestre ESP32-S3 N16R8, satélite da Equipe e satélite de Visitantes. Cada um usa duas partições de aplicação, chamadas \texttt{ota\_0} e \texttt{ota\_1}. A atualização grava a partição que não está executando e conserva a imagem anterior até a troca de boot. Essa escolha protege a operação contra download interrompido e permite ao \emph{bootloader}, o pequeno programa que escolhe o aplicativo na partida, retornar à imagem anterior se a nova falhar no primeiro ensaio. A instalação inicial \textbf{precisa ser feita por cabo} para gravar também a tabela de partições; simplesmente atualizar um firmware antigo que use uma única partição não cria o espaço OTA.

A tabela~\ref{tab:particoes_nexus} apresenta os deslocamentos e tamanhos exatos da flash do mestre e dos satélites. Cada slot de aplicação tem 6\,MiB no mestre e 1\,920\,KiB nos satélites; \texttt{otadata} mantém a seleção de partida, enquanto \texttt{nvs} conserva preferências fora das imagens. A imagem nova deve caber \emph{inteira} no slot inativo. O mestre verifica o limite antes de baixar e o satélite rejeita no comando de início um tamanho maior que o seu slot. Não se deve gravar um binário ou tabela preparados para o N16R8 em um SuperMini de 4\,MiB.

Para instalar essa tabela pela primeira vez, conecte cada placa por USB, escolha a porta correspondente no PlatformIO e grave os ambientes corretos. As instalações posteriores podem seguir por OTA, desde que as partições já estejam presentes:
\begin{lstlisting}[caption={Primeira gravação com a tabela de partições OTA}]
pio run -d firmware/mestre-s3 -e mestre-s3-n16r8 -t upload
pio run -d firmware/equipe-s3 -e esp32-s3-devkitm-1 -t upload
pio run -d firmware/visitantes-s3 -e esp32-s3-devkitm-1 -t upload
\end{lstlisting}

\section{Versões, manifesto e conjunto publicável}
A versão de cada aplicativo tem forma \texttt{MAJOR.MINOR.PATCH}. Incrementar \emph{major} anuncia incompatibilidade deliberada que requer migração ou atualização coordenada; \emph{minor} acrescenta função preservando a compatibilidade prevista; \emph{patch} corrige comportamento sem mudar o contrato. O número na etiqueta do GitHub Release identifica o \emph{conjunto} publicado. As três versões internas podem ser diferentes: uma publicação nova pode transportar um satélite alterado e dois binários inalterados. Comparar apenas a etiqueta do release produziria atualizações desnecessárias; por isso, o mestre compara separadamente a versão informada no manifesto com a versão em execução de cada alvo.

O arquivo \texttt{manifesto-firmwares.json} é o índice desse conjunto. Seu formato atual é \texttt{versao\_formato: 2}; contém \texttt{firmwares.mestre-s3}, \texttt{firmwares.equipe-s3} e \texttt{firmwares.visitantes-s3}. Para cada alvo há versão, tamanho em bytes, endereço HTTPS do binário e SHA-256. SHA-256 é um resumo criptográfico de 256 bits: uma mudança de conteúdo altera o resumo com probabilidade extremamente alta. O firmware exige endereço \texttt{https://}, resumo hexadecimal válido, versões interpretáveis e os três alvos presentes; limita o manifesto recebido a 8\,192 bytes. A versão dos satélites também precisa caber no campo de versão do protocolo UART. A etiqueta e os endereços devem corresponder ao mesmo conjunto de anexos. \textbf{HTTPS autentica o servidor por certificado; SHA-256 compara o arquivo recebido com o valor anunciado pelo manifesto.} O manifesto público, por si só, não é uma assinatura digital independente da conta e do canal de publicação; proteger permissões do repositório e do release continua necessário.

O mestre consulta \path{releases/latest/download/manifesto-firmwares.json} do repositório configurado. O release destinado à frota deve, portanto, ser o mais recente atendido por esse caminho e conter os quatro anexos completos: três \texttt{.bin} e o manifesto. O firmware atual não fornece credenciais para releases privados. Se a publicação for privada, incompleta ou apontar a binário ausente, a consulta ou a transferência falha sem instalar uma imagem parcial. O cabeçalho \texttt{Date} de uma resposta HTTPS válida pode acertar o relógio UTC dos logs no SD quando falta uma hora confiável; isso não substitui a validação TLS.

\subsection{Preparação local e publicação}
O estado de referência é \texttt{versoes/estado\_versoes.json}. Ele registra a última etiqueta, as três versões e impressões das fontes. O preparador considera código compartilhado em todos os alvos afetados e também HTML, CSS, JavaScript, geradores e páginas C geradas dos satélites. Arquivos temporários de compilação não devem determinar a versão. Depois de escolher conscientemente \texttt{patch}, \texttt{minor} ou \texttt{major}, execute na raiz do repositório:
\begin{lstlisting}[caption={Preparação local do conjunto OTA}]
python ferramentas/preparar_release_ota.py --bump minor
\end{lstlisting}
Sem \texttt{--bump}, o script oferece um menu; \texttt{--forcar equipe-s3} ou outro alvo só deve ser usado quando é necessário reconstruir e publicar nova versão sem mudança detectada. A preparação regenera painéis, detecta fontes alteradas, compila os alvos necessários e confere a versão \emph{embutida} em cada \texttt{.bin}. Imagens anteriores só são reaproveitadas quando correspondem ao estado registrado. Confirme os quatro arquivos em \texttt{versoes/saida/}, as versões internas, o tamanho e o SHA-256 de cada binário, bem como a etiqueta única indicada no manifesto. Não edite o número de versão somente no JSON: isso criaria divergência com a descrição do aplicativo gravada no binário.

Preparar arquivos locais não os publica. O caminho normal de publicação é o workflow \path{.github/workflows/publicar-firmware-mestre.yml}. Faça commit do código, do estado de versões e das páginas web geradas, envie esse commit, crie a etiqueta \texttt{vX.Y.Z} \emph{nele} e envie a etiqueta. O workflow recompila, rejeita páginas geradas desatualizadas e publica os quatro anexos no GitHub Release. Confira o resultado em Actions e Releases, inclusive as notas de migração quando houver incompatibilidade. Uma tag criada antes do commit correto, um workflow falho ou um release com anexos misturados não constitui uma versão OTA utilizável. Não reutilize uma etiqueta para conteúdo diferente, pois isso torna ambíguo o diagnóstico de qual imagem foi instalada.

\section{Consulta e autorização no equipamento}
O intervalo padrão de consulta é de 15 minutos; a configuração admite de 1 minuto a 7 dias. O mestre abre uma sessão Wi-Fi de finalidade OTA, busca o manifesto por HTTPS, compara as versões e fecha a sessão quando possível. O prazo HTTP configurado é de 20 segundos por operação. O Wi-Fi não fica ligado continuamente apenas para procurar versões. Uma sessão de outra finalidade em andamento impede uma verificação concorrente, evitando disputa pelo rádio, pela conexão e pelo estado da operação.

O satélite somente aparece como candidato se estiver respondendo recentemente, anunciar capacidade OTA UART e tiver versão inferior à do manifesto. O mestre considera a resposta de telemetria recente até 5 segundos para avaliar o enlace. A ausência de resposta, capacidade ou versão nova é mostrada como indisponibilidade, falta de suporte ou atualização desnecessária, conforme o caso. Uma oferta de atualização permanece aguardando autorização por \textbf{5 minutos}; expirado o prazo, é descartada sem alterar o firmware em execução. Uma nova consulta pode oferecer a mesma versão novamente. Esse prazo evita que uma autorização tardia instale uma decisão tomada com informação antiga sobre enlace, versão ou publicação.

No terminal serial do mestre a 115\,200 baud, a opção \texttt{2} abre o submenu OTA: \texttt{2} verifica agora; \texttt{3}, \texttt{4} e \texttt{5} autorizam, respectivamente, mestre, Equipe e Visitantes; \texttt{6} atualiza todos; \texttt{7} cancela; \texttt{8} altera o intervalo de consulta; \texttt{9} altera o intervalo de telemetria durante OTA. A página \emph{Comandos} do ponto de acesso da Equipe oferece autorizações equivalentes. Quem se associa ao AP Equipe pode enviar a ação HTTP sem uma chave administrativa adicional. A ponte Equipe--mestre devolve resposta do comando, mas o andamento deve ser consultado no estado OTA. Uma resposta de pedido aceito significa que o serviço registrou a autorização; \textbf{não significa que o download começou, que os bytes chegaram ou que o novo aplicativo foi confirmado}. A ação para \emph{todos} depende dos três alvos; se um satélite estiver ausente, pode-se selecionar apenas o alvo que responde, mantendo a verificação de versão e do estado final.

Na opção \emph{Atualizar todos}, a ordem é Equipe, Visitantes e mestre. Alvos já atualizados são pulados. Falha de um satélite interrompe a sequência antes de substituir o mestre. O motivo é operacional: o mestre antigo permanece disponível para diagnóstico e para repetir a tentativa por alvo. Uma atualização de só um alvo é adequada quando a mudança é compatível; para alteração \emph{major} do protocolo, as notas do release devem especificar a ordem segura de migração e a compatibilidade de versões intermediárias.

\section{Instalação do mestre por HTTPS}
Após a autorização, o mestre escolhe com \texttt{esp\_ota\_get\_next\_update\_partition} o slot inativo. \texttt{esp\_https\_ota} baixa e grava a imagem nesse slot, usando validação do certificado TLS pelo conjunto de raízes configurado. Antes de qualquer troca de boot, compara a versão declarada \emph{dentro} do binário com a versão prometida no manifesto, confere o tamanho informado quando disponível, rejeita imagem que exceda o slot, confere a quantidade total efetivamente recebida e calcula o SHA-256 da partição gravada. O gerenciador só chama \texttt{esp\_https\_ota\_finish} após esses testes; nessa etapa a partição passa a ser a escolhida para a próxima partida. Em caso de erro ou cancelamento antes disso, aborta a operação e a partição anterior continua executando.

O serviço registra bytes recebidos, porcentagem em décimos, taxa e estimativa de tempo restante. Quando a gravação termina com sucesso, informa estado \emph{aplicado} e reinicia após aproximadamente 3 segundos. Esse estado indica \textbf{imagem escrita, íntegra e selecionada para a primeira partida}; ainda não demonstra que as tarefas, Wi-Fi, UART e demais serviços funcionam com o novo firmware. A confirmação funcional ocorre depois do reboot e da janela de observação descrita adiante.

\section{Instalação dos satélites: HTTPS até o mestre, UART até o destino}
Os satélites não precisam acessar o GitHub. O mestre baixa por HTTPS o binário específico do alvo e o encaminha em blocos pela UART dedicada. O serviço usa três buffers reutilizáveis de 4 KiB, preferencialmente na PSRAM e, se ela não estiver disponível, tenta alocá-los na RAM interna. Uma tarefa pode preencher o próximo buffer pela rede enquanto outra transmite e aguarda confirmação do bloco anterior. Isso evita manter o firmware inteiro na RAM ou no cartão SD. O quadro entregue ao driver UART precisa de memória interna apropriada ao driver; a PSRAM serve como área de encadeamento, não como substituição irrestrita de buffers do periférico.

O mestre escolhe 4\,096 bytes por bloco quando o satélite anuncia essa capacidade. Para firmware compatível que só anuncia OTA UART, usa 2\,048 bytes. Durante o envio, a telemetria não é simplesmente desligada: o intervalo padrão passa a 2 segundos e pode ser ajustado entre 0,5 e 10 segundos. Isso conserva observabilidade básica, mas deixa mais tempo de UART para a imagem. A velocidade real depende de rede, UART, gravação da flash e confirmações; o percentual é calculado a partir dos \emph{bytes confirmados} pelo satélite, não apenas baixados pelo mestre.

\begin{table}[htbp]\centering\small
\caption{Quadro OTA UART versão 1 dos satélites}\label{tab:quadro-ota-satelite}
\begin{tabularx}{\textwidth}{@{}l l X@{}}\toprule
Campo & Tamanho & Finalidade\\\midrule
Marcador \texttt{A5 5A} e versão & 3 bytes & Identificam a família de quadro e o formato.\\
Tipo & 1 byte & Iniciar, dados, finalizar, cancelar ou resposta.\\
Sessão & 4 bytes & Distingue uma transferência de outra.\\
Sequência & 4 bytes & Ordena quadros e permite reconhecer repetição.\\
Deslocamento & 4 bytes & Início do bloco, ou próximo byte esperado na resposta.\\
Tamanho da carga & 2 bytes & Permite delimitar até 4\,096 bytes.\\
Carga & variável & Dados da imagem ou metadados e resultado.\\
CRC-32 & 4 bytes & Detecta alteração acidental do quadro recebido.\\\bottomrule
\end{tabularx}\end{table}

O cabeçalho compactado soma 18 bytes; campos multibyte são transmitidos em \emph{little-endian}. O CRC-32 verifica cada quadro, enquanto SHA-256 verifica o arquivo completo: o primeiro ajuda a decidir se um bloco deve ser repetido; o segundo impede selecionar uma imagem final diferente da anunciada. A carga \emph{INICIAR} tem 60 bytes: tamanho total de 32 bits, SHA-256 binário de 32 bytes e versão esperada em 24 bytes. A resposta contém código, fase e erro ESP-IDF; o cabeçalho devolve sessão, sequência e próximo deslocamento. As fases incluem ocioso, recebendo, validando, pronto para reiniciar, cancelado e falha. O identificador de sessão é sorteado pelo mestre; um cancelamento antigo não pode encerrar uma nova transferência.

No início, o satélite verifica formato, versão, tamanho e espaço do slot inativo, abre a gravação por \texttt{esp\_ota\_begin} e responde que está recebendo. Em cada bloco, exige sessão, sequência e deslocamento esperados antes de chamar \texttt{esp\_ota\_write}. Se a resposta se perder e o mestre repetir o último bloco já gravado, o satélite reconhece a repetição e confirma o mesmo próximo deslocamento sem gravar novamente. Sequência ou posição incoerente recebe erro explícito. O mestre valida sessão, sequência, fase e deslocamento da resposta: uma confirmação de outro bloco não avança o progresso.

Cada envio admite até quatro tentativas. O prazo de resposta é 1,5 segundo para dados, 10 segundos para iniciar e 15 segundos para finalizar. Esses prazos maiores envolvem apagar e preparar a partição ou completar a verificação final. Sem dados da sessão por 15 segundos, o satélite aborta a transferência parcial por inatividade; a verificação continua mesmo se trafegarem outros tipos de quadro na UART. Se o pipeline interno do mestre ultrapassa seu próprio prazo, o serviço o bloqueia até uma reinicialização segura em vez de arriscar reutilizar buffers ou respostas de uma operação indeterminada.

O comando \emph{FINALIZAR} só é aceito quando todos os bytes esperados foram gravados. O satélite chama \texttt{esp\_ota\_end}, calcula o SHA-256 da partição, compara com o resumo anunciado, lê a descrição do aplicativo e compara sua versão com a esperada. Só então seleciona a partição de boot, transmite a fase \emph{pronto para reiniciar} e reinicia cerca de 1 segundo depois. Ele tenta emitir essa resposta três vezes; se nenhuma transmissão for bem-sucedida, tenta restaurar a partição anterior e cancela o reinício. Se o mestre não recebe confirmação coerente, considera a tentativa falha, mesmo que alguns bytes tenham sido recebidos. Depois do reinício, espera por até 15 segundos até que uma resposta de telemetria recente anuncie a versão esperada. Esse é o critério de \emph{conclusão da transferência} no mestre; a confirmação interna de estabilidade do satélite ainda ocorrerá no prazo de 2 minutos.

\section{Estados da imagem e duas janelas de proteção}
O \texttt{otadata} registra qual slot deve iniciar e o estado da imagem. No mecanismo de rollback do ESP-IDF, a imagem recém-selecionada entra como \texttt{NEW}; na primeira partida passa a \texttt{PENDING\_VERIFY}. Se o aplicativo reiniciar antes de confirmar, o bootloader marca essa tentativa como abortada e procura a imagem válida anterior. Se os serviços essenciais falharem durante a condição pendente, o aplicativo solicita \path{esp_ota_mark_app_invalid_rollback_and_reboot}. Após iniciar e funcionar por 2 minutos, chama \path{esp_ota_mark_app_valid_cancel_rollback}, que marca a imagem válida. O \texttt{otadata} tem 8 KiB, organizados pelo ESP-IDF para conservar os dados de seleção entre gravações e reinícios; o binário continua nos slots \texttt{ota\_0}/\texttt{ota\_1}.

\begin{figure}[htbp]\centering
\begin{tikzpicture}[node distance=7mm, >={Latex},
estado/.style={draw=nexus,fill=pale,rounded corners,align=center,text width=9.6cm,minimum height=9mm,font=\small}]
\node[estado] (anterior) {Imagem anterior válida; slot inativo disponível};
\node[estado,below=of anterior] (escrita) {Download e gravação no slot inativo; tamanho, versão e SHA-256 conferidos};
\node[estado,below=of escrita] (novo) {Slot novo selecionado; bootloader inicia imagem em \texttt{PENDING\_VERIFY}};
\node[estado,below=of novo] (dois) {Serviços essenciais ativos e 2 minutos de observação};
\node[estado,below=of dois] (cinco) {Imagem \texttt{VALID}; guarda RTC acompanha reinícios anormais até 5 minutos estáveis};
\node[estado,below=of cinco] (fim) {Observação encerrada; imagem em operação normal};
\draw[->,thick] (anterior)--(escrita);
\draw[->,thick] (escrita)--(novo);
\draw[->,thick] (novo)--(dois);
\draw[->,thick] (dois)--(cinco);
\draw[->,thick] (cinco)--(fim);
\end{tikzpicture}
\caption{Caminho normal de uma imagem OTA. Falhas, cancelamento e retorno à imagem anterior estão descritos na tabela~\ref{tab:falhas-ota}.}\label{fig:estados-ota}
\end{figure}

O mestre e os dois satélites registram a partida antes de completar os serviços. Durante o estado pendente, falha considerada essencial no início dos serviços provoca rejeição da imagem e pedido de retorno; em especial, falhas de Wi-Fi ou do serviço OTA comprometem a capacidade de recuperação. Depois de inicializar os serviços, a aplicação aguarda 120\,000 ms, confirma a imagem em execução e mantém observação até 300\,000 ms de estabilidade. O primeiro prazo dá tempo para defeitos de inicialização e operação imediata aparecerem \emph{antes} de remover a proteção do bootloader. O prazo total de 5 minutos estende a detecção para falhas que dependem de alguma atividade ou tempo de execução.

Após a confirmação de 2 minutos, o projeto usa ainda uma \emph{guarda} em memória RTC não inicializada (\texttt{RTC\_NOINIT\_ATTR}). Ela registra assinatura, endereço da partição em execução, contador de reinícios anormais e um valor de integridade calculado por XOR. Na partida, só aceita o registro se assinatura, integridade e endereço da partição coincidirem. Essa checagem evita interpretar resíduos de memória como estado válido ou atribuir a contagem a outro firmware. O estado não requer gravações frequentes na flash, preservando sua vida útil. A memória RTC costuma sobreviver a resets de software e de watchdog, mas \textbf{não é garantia diante de corte total de energia}; portanto, a contagem não deve ser apresentada como histórico persistente de falhas.

A guarda classifica como reinício anormal \texttt{ESP\_RST\_PANIC}, \texttt{ESP\_RST\_INT\_WDT}, \texttt{ESP\_RST\_TASK\_WDT} e \texttt{ESP\_RST\_WDT}. Após a confirmação inicial, cada partida por um desses motivos incrementa o contador. Ao atingir três, chama a API de invalidação e rollback. Reinícios intencionais ou de outra causa não entram nessa contagem. Ao completar o período estável de 5 minutos, a aplicação apaga o estado RTC da guarda. A guarda complementar não substitui o primeiro ensaio do bootloader, e o retorno depende de existir uma partição anterior válida. Se ela não existir, a API pode falhar; nesse caso, a recuperação exige diagnóstico e possível gravação por cabo.

\begin{nota}{Três significados de confirmação}
\textbf{Pedido aceito:} a interface HTTP ou o comando UART aceitou iniciar uma operação. \textbf{Imagem transferida:} os bytes foram gravados, o resumo e a versão conferem; no satélite, o mestre também viu a nova versão responder após o reboot. \textbf{Imagem confirmada como válida:} o aplicativo novo executou seus serviços essenciais por 2 minutos e chamou a confirmação do ESP-IDF. O acompanhamento de reinícios continua até 5 minutos. A palavra \emph{validada} em um log logo após a gravação refere-se à integridade do arquivo, não à aprovação funcional de 2 minutos.
\end{nota}

\section{Cancelamento, falhas e recuperação prática}
Antes de autorizar, \emph{Cancelar} descarta a oferta. Durante um download do mestre, a solicitação é verificada entre etapas de \texttt{esp\_https\_ota\_perform}; o gerenciador aborta antes de selecionar o slot novo. Durante transferência a satélite, o mestre para o pipeline e envia \emph{CANCELAR} com a sessão ativa; o satélite aborta o manipulador OTA e mantém o firmware atual. Um \emph{CANCELAR} de outra sessão é ignorado. Cancelar depois de o destino ter escolhido uma imagem e iniciado o reboot não é equivalente a desfazer uma instalação: nesse ponto, observar a versão que voltou a responder e o estado de rollback. A ação pode ser aceita antes de a tarefa encerrar; acompanhe o estado final e o erro em vez de repetir imediatamente comandos.

\begin{table}[htbp]\centering\small
\caption{Interpretação de falhas durante OTA}\label{tab:falhas-ota}
\begin{tabularx}{\textwidth}{@{}l X X@{}}\toprule
Indício & Causa provável a investigar & Primeira verificação\\\midrule
Manifesto ausente ou inválido & Release incompleto, URL inacessível, formato divergente ou Wi-Fi sem Internet. & Conferir sessão Wi-Fi, resposta HTTPS, \texttt{versao\_formato} e quatro anexos do release.\\
Alvo indisponível ou sem suporte & Telemetria antiga, UART desconectada, firmware satélite sem capacidade OTA. & Conferir GND comum, pinos RX/TX, versão e idade da última resposta.\\
Bloco repetido ou timeout & Perda de resposta UART, ruído, excesso de tráfego ou satélite ocupado gravando. & Conferir sequência, deslocamento, CRC, contadores de erro e cabos; não interpretar repetição como dupla gravação.\\
Imagem rejeitada ao finalizar & Tamanho, SHA-256, versão interna ou formato de aplicativo divergente. & Comparar binário e manifesto do mesmo release; recompilar o conjunto se necessário.\\
Versão nova não responde & Partida falhou, enlace não voltou ou rollback ocorreu. & Aguardar o prazo, ler logs seriais dos dois dispositivos e consultar versão real antes de repetir.\\
Volta à versão anterior & Falha no primeiro ensaio, serviço essencial ou três resets anormais na janela RTC. & Ler motivo de reset e erro de serviço; corrigir e publicar nova etiqueta.\\\bottomrule
\end{tabularx}\end{table}

Para instalar, comece com as três placas energizadas, GND comum e versões conhecidas. Confirme no painel da Equipe que mestre e satélites respondem. No terminal, abra OTA, verifique agora e anote versões instalada e oferecida. Autorize o alvo ou \emph{todos} dentro de 5 minutos. Observe bytes confirmados, taxa, erro e versão reportada após cada reinício; o terminal e a OLED do mestre apresentam o mesmo estado, incluindo uma barra de progresso. Aviso sonoro distingue atualização disponível, conclusão e falha. Aguarde a janela de 2 minutos para confirmação funcional e a de 5 minutos para encerramento da guarda antes de declarar a instalação estável. Se a oferta expirar, faça nova consulta; se um satélite estiver ausente, recupere o enlace antes de repetir.

O fluxo foi implementado com APIs OTA do ESP-IDF. Para detalhes do significado de \texttt{NEW}, \texttt{PENDING\_VERIFY}, \texttt{VALID}, \texttt{INVALID} e \texttt{ABORTED}, consulte a \href{https://docs.espressif.com/projects/esp-idf/en/latest/esp32s3/api-reference/system/ota.html}{documentação oficial de OTA do ESP-IDF}. Os procedimentos específicos do NEXUS, parâmetros de release e comandos estão registrados em \texttt{documentacao/ota.md}; os estados e prazos descritos aqui correspondem aos serviços e à guarda de rollback do repositório.

\chapter{Diagnóstico, supervisão e apresentação local}\label{cap:diagnostico}
\section{Diagnóstico e causas de falha}
Um valor aparentemente plausível pode estar antigo, vir de um sensor sem calibração ou sobreviver na tela após a perda de um cabo. Por isso, o mestre acompanha a origem e a idade dos dados, erros de drivers, respostas UART, memória, tarefas e falhas de armazenamento. O diagnóstico junta fatos coletados pelos módulos e publica estado \emph{saudável}, \emph{atenção}, \emph{falha} ou \emph{recuperando}. A interface web, o terminal, os LEDs e a OLED leem esse retrato em vez de consultar drivers diretamente. Contadores históricos permanecem úteis mesmo quando uma causa ativa desaparece.

\subsection{Falhas que o mestre identifica}
O ciclo de saúde é atualizado a cada segundo. Os serviços de MPU, encoders e INA228 informam seus resultados; três erros consecutivos de um serviço de sensor abrem uma causa de atenção. O gerenciador I\textsuperscript{2}C conta transações bem-sucedidas, falhas e endereços ausentes. Uma rajada de três erros em aproximadamente 15\,s mantém aviso I\textsuperscript{2}C por dez segundos. UART mantém erros de transmissão e recepção e descartes; rajada de três erros ou novo descarte produz aviso por cinco segundos. Falha OTA recente aparece por dez segundos. Os contadores históricos permanecem disponíveis após a causa ativa desaparecer.

Uma confirmação de satélite é considerada recente quando sua idade não supera três vezes o intervalo atual de telemetria mais 500\,ms. Depois de ao menos dez pacotes gerados e três ciclos de diagnóstico sem resposta, surge a causa Equipe ou Visitantes sem resposta. Isso evita alarme durante a primeira partida e adapta o prazo ao intervalo mais lento de OTA. RTT atual, idade da última confirmação, perdas, duplicatas e versões ajudam a diferenciar satélite desligado, fio interrompido, parser incompatível e congestionamento.

O mestre mede sinal da rede externa apenas quando conectado. RSSI abaixo de --85\,dBm por dez amostras gera aviso; acima de --80\,dBm por cinco amostras limpa o aviso. A faixa intermediária produz histerese para evitar alternância perto do limiar. RSSI é potência recebida aproximada, não taxa de dados nem prova de acesso à Internet. A telemetria UART independe dessa conexão.

\begin{table}[htbp]\centering\small
\caption{Causas de saúde e observações relacionadas}\label{tab:causas-saude}
\begin{tabularx}{\textwidth}{@{}l X X@{}}\toprule
Causa & Detecção no código & Verificação útil\\\midrule
Sensor & Erros consecutivos de MPU ou encoder. & Alimentação, endereço, cabo, eixo, contador e último erro.\\
I\textsuperscript{2}C & Rajada de erros no barramento. & Pull-ups, nível alto, cabos, endereços e forma de onda.\\
UART & Erros recentes ou novo descarte. & TX/RX/GND, baud, versão, CRC e fila.\\
Satélite sem resposta & ACK ausente após início do fluxo. & Energia do satélite, idade, RTT e UART exclusiva.\\
Supervisão & Tarefa atrasada ou em recuperação. & Fase, último pulso, pilha, fila e serviço bloqueante.\\
Memória/pilha & Limiar de RAM, PSRAM, maior bloco ou pilha. & Uso sob TLS/OTA/SD e maior alocação exigida.\\
Wi-Fi fraco & RSSI persistentemente baixo na sessão. & Distância e obstáculos; manter registro do canal.\\
OTA recente & Falha da última operação. & Alvo, etapa, erro, versão que voltou e manifesto.\\
Persistência & Falha ao guardar odômetro/estado. & NVS, cartão e contadores; validar a classificação ativa em bancada.\\\bottomrule
\end{tabularx}\end{table}

\subsection{RAM, maior bloco livre e margem de pilha}
O mestre mede RAM interna e PSRAM livres, além do maior bloco contíguo alocável. Uma soma alta de bytes livres não garante que exista um bloco grande para TLS, buffer ou driver. A atenção por RAM interna exige 15 amostras com menos de 24\,KiB livres ou maior bloco abaixo de 8\,KiB; falha exige três amostras abaixo de 16\,KiB ou 4\,KiB. Abaixo de 8\,KiB livres ou 2\,KiB de bloco, a falha é imediata. Para PSRAM presente, os limiares correspondentes são atenção abaixo de 256\,KiB livres ou bloco de 32\,KiB; falha abaixo de 128\,KiB ou bloco de 16\,KiB; emergência abaixo de 64\,KiB ou bloco de 8\,KiB. Cinco amostras fora da faixa limpam o alerta. Essa persistência evita alarmes causados por uso temporário legítimo durante HTTPS.

A margem mínima de pilha de tarefa do aplicativo abaixo de 256 bytes é causa crítica. O sistema coleta uma tabela de tarefas a cada cinco segundos, com prioridade, afinidade, estado, uso de CPU por janela e menor margem observada. A primeira janela de CPU não possui base para comparação. Uma CPU de 100\% representa um núcleo; afinidade não informa em qual núcleo a tarefa executou naquele instante. Tarefas do SDK aparecem para diagnóstico, mas a classificação de pilha crítica usa a política das tarefas do projeto. As pilhas do \texttt{main} e da tarefa de eventos do ESP-IDF foram ajustadas na configuração para 6\,144 e 4\,096 bytes após margens pequenas observadas em bancada; alterações futuras devem se apoiar em novas medições.

\subsection{Supervisão e recuperação}
Serviços alimentam pulsos de progresso e informam fase e prazo. Uma tarefa de alta prioridade detecta atraso e enfileira recuperação; uma tarefa separada executa ações que podem bloquear. Cada UART tem recuperação independente de filas, buffers, pinos e velocidade. Para telemetria atrasada, o supervisor antecipa o próximo ciclo por notificação, sem apagar à força uma tarefa que pode estar usando um driver. SD, aquisição, upload e OTA também têm fases observáveis; a espera ociosa usa prazo zero. O objetivo é tornar travamentos localizáveis e preservar os módulos que continuam funcionais.

\section{O que a OLED mostra}
A SSD1306 de 128$\times$64 pixels usa o mesmo I\textsuperscript{2}C do MPU. A tela reserva até oito linhas curtas; quando há progresso, a última faixa vira barra gráfica. Na partida, apresenta etapa e percentual (NVS, supervisão, sensores/enlaces, cartão, telemetria, rede, OTA, terminal e diagnóstico), inclusive aviso de etapa degradada. Depois, em operação saudável, alterna aproximadamente a cada seis segundos entre três páginas: estado/tempo ativo/RAM/PSRAM/versão; telemetria (velocidade, bateria, tensão, potência, marcha, sequência e origem real ou simulada); e comunicação (resposta Equipe/Visitantes, RTT e Wi-Fi).

\begin{nota}{Limite importante da página de telemetria local}
O estado de carga da bateria não decorre automaticamente da leitura do INA228. A tensão e a potência vêm do sensor apenas quando a medição é válida; módulo ausente, erro de barramento ou simulador mudam a origem e a disponibilidade. A OLED tem pouco espaço para sinalizar validade individual e pode exibir um zero de campo não medido: \textbf{zero isolado não comprova bateria vazia, tensão zero ou potência zero}. Conferir origem e validade no painel técnico ou no terminal antes de usar a tela para uma decisão operacional.
\end{nota}

Durante uma sessão Wi-Fi, a OLED mostra rede ativa, RSSI, canal, tentativas, prazos e último erro; quando o rádio está parado, a página normal informa repouso. Na verificação OTA, mostra fase de conexão e consulta HTTPS, rede e versão. Com atualização disponível, lista versões de mestre, equipe e visitantes e a necessidade de autorização. Na transferência, mostra alvo, versão, bytes/porcentagem, taxa e tempo restante; durante upload ao Google, arquivo atual, progresso, velocidade, prazo e confirmação. Após upload, apresenta resumo por cerca de oito segundos. Avisos alternam causa a cada três segundos quando há várias e recomendam consultar o terminal.

Uma fila fixa guarda mensagens curtas de sucesso, atenção ou erro. Eventos importantes reservam a tela; mensagens adiadas são mostradas depois, sem alocação dinâmica por aviso. Atualização disponível e falhas reaparecem em janelas, mas devolvem a tela às páginas normais entre lembretes. A tarefa só envia novo quadro ao SSD1306 quando o conteúdo muda; falha de atualização é contada e tentada novamente. Ausência da OLED não interrompe aquisição nem UART.

\section{LEDs, buzzer, logs e onde investigar}
Quatro LEDs e um buzzer passivo dão sinais mesmo sem navegador. Verde contínuo representa saúde; amarelo periódico, atenção; vermelho quase contínuo com breve pausa, falha crítica. Azul indica partida, Wi-Fi ou OTA por combinações distintas. O serviço escolhe prioridade entre saúde, operação, Google e atualização; a máquina de tons avança a cada 50\,ms sem travar a aquisição. Sons marcam transições, como sistema pronto, oferta OTA, conclusão ou erro. Wi-Fi de rotina e verificações periódicas silenciosas evitam alarmes repetitivos.

O terminal serial traz o detalhe que não cabe na OLED. O menu de sensores permite pausar todas as leituras ou fontes individuais para depuração, inclusive o INA228; consultar o estado de pausa antes de atribuir ausência de dado a uma falha elétrica. O menu de logs ajusta o nível global e por serviço, incluindo a opção de silenciar avisos no terminal e de ligar ou desligar cores ANSI. Essas preferências de apresentação não apagam os contadores de falha nem substituem os registros do SD. O painel técnico consulta retratos compactos de CPU, memória, temperatura, tarefas e causas do mestre, além da saúde local da Equipe e da resposta de Visitantes. O SD guarda logs estruturados e uma janela de histórico pré-falha em PSRAM quando disponível; sem cartão, essa janela pode permanecer apenas em memória volátil. Em diagnóstico, anotar hora relativa ao boot, versão, estado dos três enlaces, últimos erros, causas ativas e arquivos brutos. Uma causa que desapareceu pode continuar no contador histórico, e um contador crescente não significa necessariamente falha presente.

\chapter{Instalação, operação, testes e manutenção}\label{cap:operacao}
\section{Preparação do ambiente e primeira partida}
Instalar Python e PlatformIO com suporte a ESP-IDF, obter a revisão correta do repositório e registrar o commit. Conferir \texttt{platformio.ini} de cada alvo e as partições OTA. Os painéis web são gerados a partir das fontes HTML/CSS/JavaScript por \texttt{gerar\_painel.py}; o arquivo C incorporado deve corresponder a elas. Credenciais Wi-Fi, URL e token do Google pertencem à configuração local ignorada pelo Git. Não colocar segredos no ZIP do Overleaf nem no histórico público.

Na primeira instalação, gravar cada placa por cabo para criar a tabela de partições. Os comandos e ambientes corretos estão no capítulo~\ref{cap:ota}; selecionar a porta e identificar fisicamente cada placa antes da gravação.
O ambiente exato do mestre fixa \texttt{espressif32@6.13.0}; os satélites devem ser compilados com dependências registradas no ensaio, pois seus arquivos não fixam essa versão. O terminal serial do mestre e a chave inicial da equipe usam 115\,200 baud. A gravação por cabo não dispensa conferir GPIO, alimentação e conexão física antes de ligar o conjunto.

\subsection{Montagem e partida}
\begin{enumerate}
\item Identificar a revisão física de cada PCB e o esquema correspondente; verificar polaridade, tensão dos conversores, proteções, continuidade de GND e ausência de curto com a alimentação desligada.
\item Confirmar 3,3\,V nos domínios lógicos. Conferir pull-ups dos quatro canais de encoder para 3,3\,V, nunca para a alimentação externa. Conferir TX cruzado com RX em cada UART.
\item Confirmar geometria das rodas e relação encoder/roda; conferir orientação X/Y do MPU; registrar valores efetivos antes de análise quantitativa.
\item Energizar, observar serial, LEDs e OLED. Verificar se o mestre inicia, se o sensor inercial calibra em repouso, se os encoders não contam com rodas paradas e se o cartão monta quando instalado.
\item Acessar o AP da Equipe em \url{http://192.168.4.1/}; conferir versão e idade da telemetria dos dois satélites, sequência e ACK. Verificar se a origem indicada é real ou simulada.
\item Iniciar um registro curto, conferir o contador de amostras no SD e preservar o CSV original. Para ensaio do painel público, verificar que Visitantes mostra somente leitura.
\end{enumerate}

\subsection{Conferência elétrica antes da primeira energização}\label{sec:conferencia-eletrica}
Trabalhar com a alimentação de tração isolada e descarregada segundo o procedimento elétrico da equipe. Identificar a revisão efetivamente montada; não usar apenas os nomes dos conectores da folha v4. Com multímetro, confirmar polaridade dos trilhos, continuidade dos retornos, ausência de curto entre 48 V, 5 V, 3,3 V e GND, e orientação do diodo de roda livre. Energizar a eletrônica por fonte limitada em corrente antes de conectá-la ao veículo. Medir 5 V e 3,3 V em cada módulo durante partida, Wi-Fi e escrita no cartão, além da temperatura do buck e do regulador. Registrar queda de tensão e ruído com osciloscópio quando possível. O capítulo \ref{cap:eletronica} explica as razões elétricas e os pontos não comprovados do esquema.

Para o I\textsuperscript{2}C, contar todos os resistores de pull-up soldados na placa e nos módulos. Medir SDA--3,3 V e SCL--3,3 V com o conjunto desligado, comparar a resistência equivalente com capacitância e tempo de subida, e retirar os pull-ups excedentes dos módulos quando necessário. Conferir que nenhum módulo leva as linhas a 5 V. Observar uma transação a 400 kHz com todos os módulos conectados; registrar nível alto, nível baixo, subida e ausência de colisões. Só então verificar endereços de MPU, OLED e INA228, caso este último esteja instalado. O driver do INA228 está implementado, mas sua resposta depende de módulo presente, ligação correta e configuração coerente com o shunt físico.

Para o shunt, confirmar que o componente real é de 1 m$\Omega$, sua potência admissível contínua e a montagem térmica a 75 A. Usar ligação Kelvin independente dos cabos de potência, marcar IN+ e IN$-$ e comparar a tensão diferencial com $1\,\mathrm{mV/A}$ em correntes conhecidas e seguras. Medir aquecimento em regime; 75 A produzem 5,625 W no elemento resistivo. Não usar leitura de corrente, potência ou estado de bateria como proteção até que os limites e os ensaios da montagem sejam aceitos; o driver de leitura, sozinho, não implementa proteção de tração.

O relé e seu transistor devem ser ensaiados separados do controlador de tração: identificar tensão nominal da bobina, corrente com o contato aberto e fechado, estado em reset e falta de alimentação, polaridade de D1, tempo de soltura e contato normalmente aberto/fechado. A topologia representada na folha v4 exige revisão antes de qualquer conexão ao driver. O estado remoto de parada descrito no capítulo \ref{cap:comandos} não aciona esse relé.

\subsection{Configuração de rede, painéis e registro}\label{sec:config-operacional}
Depois de gravar os três ESPs, observar no serial se as partições e a NVS foram inicializadas sem erro. Conectar o computador ou telefone ao AP da Equipe e abrir \url{http://192.168.4.1/}; registrar nome da rede, canal, versão, número de clientes e idade da telemetria. O acesso ao AP Equipe permite enviar comandos sem chave HTTP adicional; reservar a senha desse AP à equipe. Conferir a página de Visitantes separadamente para garantir que ela apresenta dados e não oferece comandos. O capítulo \ref{cap:wifi-paineis} descreve as rotas, a política de clientes e o motivo de o mestre ligar o Wi-Fi externo somente durante sessões.

Na Equipe ou no terminal do mestre, inspecionar perfis Wi-Fi salvos, tentar a rede principal e registrar canal, RSSI, tentativas, erro e prazo. Para uma rede secundária, registrar autorização e a política de conexão automática antes de testar. Uma conexão ao roteador não comprova acesso ao Google nem integridade do relógio ou da URL HTTPS. Fazer uma consulta OTA e um envio de CSV de ensaio separadamente, observando o estado final de cada tarefa. Para o Google, seguir o cadastro de propriedades, implantação e teste do Apps Script do capítulo \ref{cap:memorias_google}; não copiar token, chave ou SSID para este manual.

\subsection{Ensaios de falha e recuperação}\label{sec:ensaios-falha}
Verificar deliberadamente, em bancada, situações que o diagnóstico deve distinguir: sensor I\textsuperscript{2}C ausente, pull-up inadequado, UART de um satélite interrompida, Wi-Fi externo indisponível, cartão SD removido antes da partida, memória sob carga de TLS, resposta HTTP de upload inválida e imagem OTA de teste que falhe a validação. Para cada caso, registrar causa ativa, contador, página OLED/LED, resposta web, log, tempo até detecção, comportamento preservado e recuperação. Repetir após religar a alimentação; a memória RTC usada pelo guardião de OTA não é persistência garantida sob corte total. Não provocar falha por desconexão de cabos energizados de potência ou por remoção do cartão durante escrita.

Ao testar rollback, gravar uma imagem de laboratório que reinicie antes de ser validada e confirmar qual partição foi selecionada no próximo boot. Em teste separado, provocar três reinícios anormais da mesma imagem após a confirmação inicial e dentro da janela de estabilização; observar contagem e retorno. Esses são dois mecanismos diferentes, descritos no capítulo \ref{cap:ota}; registrar sequência temporal, causa de reset e versão efetiva em cada reinício. Não realizar o ensaio com controle de tração conectado.

\section{Operação e manutenção}
Antes de mover o veículo, conferir alimentação, conectores, fixação dos sensores, cabos afastados da potência, estado dos LEDs e integridade da parada física do veículo conforme o procedimento próprio da equipe. O painel remoto ainda não é mecanismo de parada física. Deixar o MPU em repouso durante a calibração inicial. Confirmar validade, origem e idade da telemetria. Ao encerrar, guardar os CSVs e logs do ensaio junto da versão de firmware e das condições experimentais. Desligar conforme a sequência elétrica adotada na montagem, evitando remover o cartão durante escrita.

\subsection{Plano de verificação}
Os critérios quantitativos de aceitação devem ser aprovados pela equipe para a finalidade do veículo; não há no material consultado uma campanha de calibração conclusiva. A tabela descreve ensaios reproduzíveis e a evidência a guardar.
\begin{table}[htbp]\centering\small
\caption{Ensaios recomendados antes de declarar uma medição validada}\label{tab:ensaios}
\begin{tabularx}{\textwidth}{@{}l X X@{}}\toprule
Objeto & Procedimento & Evidência e decisão\\\midrule
Alimentação & Medir trilhos em vazio, partida e carga; inspecionar ruído e aquecimento. & Tensão, instrumento, corrente, temperatura e comparação com limites dos componentes.\\
Encoders & Girar cada roda número conhecido de voltas, nos dois sentidos e com carro apoiado; comparar distância externa. & Contagens, $D$, $i$, erro e repetibilidade por lado; confirmar ausência de sentido na versão atual.\\
MPU6050 & Calibrar parado; verificar eixos por inclinação e movimento controlado. & Estado de calibração, orientação, ruído e resposta longitudinal.\\
UART & Interromper e restaurar cada enlace; injetar quadro corrompido em bancada. & Contadores CRC/perdas/duplicatas, idade da amostra e recuperação independente.\\
SD & Gravar por tempo suficiente, retirar alimentação de modo controlado e reiniciar. & CSV legível, contadores, rotação, perda máxima observada e preservação dos arquivos.\\
OTA & Instalar versão de teste, verificar alvo, reinício, versão nova e rollback controlado. & Manifesto, SHA-256, logs, estado após observação e resultado por dispositivo.\\\bottomrule
\end{tabularx}\end{table}

Não classificar o ensaio apenas como aprovado ou reprovado. Registrar data, responsável, revisão da placa, commit, versões dos três firmwares, instrumentos, condições, dados brutos, diferença em relação à referência e critério previamente definido. Testar também o sistema com Wi-Fi externo indisponível, SD ausente e navegador desconectado: a aquisição e a telemetria UART devem continuar quando a falha for degradável.

\subsection{Atualização de firmware}
Para OTA, conferir se há release completo com manifesto e três binários coerentes. Consultar a versão disponível no menu 2 do mestre ou na página Comandos da Equipe, autorizar o alvo, acompanhar progresso e aguardar a versão nova responder e completar a observação. Uma oferta sem autorização expira; verificar novamente se necessário. Em atualização de todos, satélites vêm antes do mestre. Não interpretar confirmação HTTP como instalação concluída.

\begin{table}[htbp]\centering\small
\caption{Sintomas e primeiras verificações}\label{tab:diagnostico}
\begin{tabularx}{\textwidth}{@{}l X X@{}}\toprule
Sintoma & Verificar & Ação inicial\\\midrule
Sem amostra recente & Alimentação dos três ESPs, TX/RX/GND, baud, CRC, sequência e ACK. & Restaurar o enlace, consultar contadores e não confiar no último valor exibido.\\
Contagem com roda parada & Pull-ups a 3,3\,V, blindagem, ruído de potência e filtro PCNT. & Isolar o chicote e medir formas de onda antes de mudar a geometria.\\
Velocidade incoerente & Diâmetro, relação de transmissão, P/R e quatro bordas contadas. & Comparar contagens por volta e aferir a distância com referência.\\
MPU inválido & SDA/SCL, endereço, alimentação, eixo e condição de repouso. & Conferir I\textsuperscript{2}C e repetir calibração após fixar a orientação.\\
Sem CSV & Cartão, 3,3\,V, SPI2, montagem FAT, fila e espaço livre. & Preservar o cartão e consultar logs; telemetria ao vivo pode seguir.\\
Ação web negada & Permissão do Mestre (quando exigida), carga inválida ou idade da UART. & Conferir a mensagem de rejeição, a permissão aplicável e consultar resposta e estado separadamente.\\
OTA interrompida & Rede externa, release, manifesto, hash, espaço da partição e ACK UART. & Ler erro por alvo e versão que responde antes de repetir.\\\bottomrule
\end{tabularx}\end{table}

\subsection{Pendências técnicas rastreáveis}
\begin{enumerate}
\item Associar a placa física instalada a um projeto KiCad, revisão de esquema, PCB, conectores e mapa de alimentação; acrescentar fotos de ambas as faces e do chicote.
\item Medir diâmetro efetivo das rodas, relação real dos encoders, resolução e erro; validar orientação do MPU6050 no chassi.
\item Integrar e calibrar o INA228 e o shunt para corrente e tensão do barramento; implementar os demais sensores físicos antes de publicar seus campos como medições.
\item Registrar testes de desempenho com SD, Wi-Fi e OTA concorrentes e os prazos de detecção de amostra antiga; aprovar limites de erro e critérios por aplicação.
\item Se houver futura atuação física, projetar uma cadeia independente de segurança, intertravamentos, autenticação de operador, prazo de validade dos pedidos e retorno de estado medido. O estado preparatório atual não atende essa função.
\end{enumerate}

\appendix
\chapter{Fontes e rastreabilidade}
\section{Fontes internas utilizadas}
O conteúdo foi conferido principalmente em \path{README.md}, \path{documentacao/arquitetura.md}, \path{documentacao/filosofia.md}, \path{documentacao/mestre.md}, \path{documentacao/satelites.md}, \path{documentacao/pinagem.md}, \path{documentacao/hardware.md}, \path{documentacao/protocolos.md}, \path{documentacao/ota.md}, \path{documentacao/google.md}, \path{documentacao/controle-remoto.md} e \path{documentacao/conceitos.md}. A referência executável é o código em \path{firmware/mestre-s3}, \path{firmware/equipe-s3}, \path{firmware/visitantes-s3} e \path{firmware/compartilhado}. Os esquemas e placas candidatos estão em \path{hardware/kicad/KiCad_pcb_telemetria}. Para a seção elétrica, foram consultadas a folha de dados do INA228 e a nota sobre pull-ups I\textsuperscript{2}C da Texas Instruments, além das diretrizes de hardware do ESP32-S3 da Espressif. O valor de shunt de 1 m$\Omega$ e a capacidade de 75 A foram informados pelo projeto e precisam ser confirmados no componente instalado.
\section{Registro de configuração para cada ensaio}
Em todo ensaio, anotar: data e responsável; identificação impressa da placa e do chicote; revisão do esquema e da PCB; commit do repositório; versões de cada firmware; configuração de sensores e rodas; instrumentos de referência e sua calibração; condições de alimentação; arquivos brutos do microSD; resultado observado e eventual incidente. O repositório consultado não associa uma revisão física única à placa instalada nem contém um relatório de validação geral. Essas informações devem acompanhar a próxima revisão deste manual.
\end{document}
